% claim.tex -- AI4Math claim of record (AI-generated, 2026-10-01; instantiated
% from harness/templates/claim/claim.tex per SOP 08b).
% Machine-readable theorem anchors are the "% LEAN:" lines; scripts/paper-lint.sh
% and the claim audit check them. All Lean listings are verbatim, byte-identical
% excerpts of frozen artifacts of tasks/20260929-l14c-graphside-pipeline:
%   statement  : formalized/statement.lean (124 lines, LF, sha256 d186ff16...)
%   proof file : attempts/phase2/r3/bridge.lean (1857 lines, CRLF on the working
%                disk, sha256 399ee81f...)
% Line ranges used: statement lines 36-98 as an ordered subsequence with exactly
% three lines excluded (declared below the listing); bridge reproduced in full
% (whole file, spliced LF-normalised per the chorded/ladder/hosoya precedent;
% byte-exact original shipped in proofs/). Verification: the package-local copy
% of inspect-tex.py under audit/ (--verify; record in audit/listing-verify.txt).
\documentclass[11pt]{article}

\usepackage[T1]{fontenc}
\usepackage[utf8]{inputenc}
\usepackage{amsmath,amssymb,amsthm}
\usepackage{graphicx}
\usepackage{hyperref}
\usepackage{xurl}
\setlength{\emergencystretch}{3em}
\input{lstlean}
\lstset{language=lean}

% --- local rendering patch (2026-10-01; lineage papers/cyl3-mod23 + the
% mossad42-01-half3 glyph-map list) --- tectonic/XeTeX does not apply the
% listings literate table to multi-byte characters, so every non-ASCII codepoint
% that occurs inside THIS note's listings is mapped to a math glyph as an active
% character via the standard \lccode/\lowercase idiom (ASCII-only source, no
% extra packages). Deliberately NOT mapped: U+00B7, U+2013, U+2014 and all CJK
% ideographs / CJK punctuation (3001, 3002, FF08-FF1B) --- they route through
% xeCJK + SimSun (xeCJK sets up its own punctuation table at load time; making
% these active first aborts the run, observed 2026-09-27/09-30 in the pendant
% and mossad42 notes). Rendering only: listing source bytes are untouched, so
% the byte-exactness checks are unaffected, and the \ifdefined guard leaves the
% pdflatex path identical to the template.
\ifdefined\XeTeXversion
  {\lccode`\~="00A7 \lowercase{\global\catcode`~=\active \global\def~{\S}}}
  {\lccode`\~="00AC \lowercase{\global\catcode`~=\active \global\def~{\ensuremath{\neg}}}}
  {\lccode`\~="03B1 \lowercase{\global\catcode`~=\active \global\def~{\ensuremath{\alpha}}}}
  {\lccode`\~="03B2 \lowercase{\global\catcode`~=\active \global\def~{\ensuremath{\beta}}}}
  {\lccode`\~="2115 \lowercase{\global\catcode`~=\active \global\def~{\ensuremath{\mathbb{N}}}}}
  {\lccode`\~="2124 \lowercase{\global\catcode`~=\active \global\def~{\ensuremath{\mathbb{Z}}}}}
  {\lccode`\~="2190 \lowercase{\global\catcode`~=\active \global\def~{\ensuremath{\leftarrow}}}}
  {\lccode`\~="2191 \lowercase{\global\catcode`~=\active \global\def~{\ensuremath{\uparrow}}}}
  {\lccode`\~="2192 \lowercase{\global\catcode`~=\active \global\def~{\ensuremath{\rightarrow}}}}
  {\lccode`\~="2194 \lowercase{\global\catcode`~=\active \global\def~{\ensuremath{\leftrightarrow}}}}
  {\lccode`\~="21A6 \lowercase{\global\catcode`~=\active \global\def~{\ensuremath{\mapsto}}}}
  {\lccode`\~="2200 \lowercase{\global\catcode`~=\active \global\def~{\ensuremath{\forall}}}}
  {\lccode`\~="2203 \lowercase{\global\catcode`~=\active \global\def~{\ensuremath{\exists}}}}
  {\lccode`\~="2208 \lowercase{\global\catcode`~=\active \global\def~{\ensuremath{\in}}}}
  {\lccode`\~="2209 \lowercase{\global\catcode`~=\active \global\def~{\ensuremath{\notin}}}}
  {\lccode`\~="2212 \lowercase{\global\catcode`~=\active \global\def~{\ensuremath{-}}}}
  {\lccode`\~="2227 \lowercase{\global\catcode`~=\active \global\def~{\ensuremath{\wedge}}}}
  {\lccode`\~="2228 \lowercase{\global\catcode`~=\active \global\def~{\ensuremath{\vee}}}}
  {\lccode`\~="222A \lowercase{\global\catcode`~=\active \global\def~{\ensuremath{\cup}}}}
  {\lccode`\~="2260 \lowercase{\global\catcode`~=\active \global\def~{\ensuremath{\neq}}}}
  {\lccode`\~="2264 \lowercase{\global\catcode`~=\active \global\def~{\ensuremath{\leq}}}}
  {\lccode`\~="2265 \lowercase{\global\catcode`~=\active \global\def~{\ensuremath{\geq}}}}
  {\lccode`\~="22A2 \lowercase{\global\catcode`~=\active \global\def~{\ensuremath{\vdash}}}}
  {\lccode`\~="25A1 \lowercase{\global\catcode`~=\active \global\def~{\ensuremath{\square}}}}
  {\lccode`\~="27E8 \lowercase{\global\catcode`~=\active \global\def~{\ensuremath{\langle}}}}
  {\lccode`\~="27E9 \lowercase{\global\catcode`~=\active \global\def~{\ensuremath{\rangle}}}}
\fi
\ifdefined\XeTeXversion
  \usepackage{xeCJK}
  \setCJKmainfont{SimSun}
\fi

\newtheorem{theorem}{Theorem}

\title{Matchings of the $n$-prism with one spoke or one rim edge deleted:\\
two bridge theorems to recurrence-defined sequences\\
{\large (Lean 4 machine-checked claim of record)}}
\author{Aurora0134\thanks{Contact: via the GitHub account Aurora0134.}}
\date{October 1, 2026}

\begin{document}
\maketitle

\begin{abstract}
Let $\mathrm{aspoke}(n)$ and $\mathrm{arim}(n)$ be the integer sequences with
initial values $25, 86, 271, 876$ and $26, 86, 274, 883$ respectively, both
satisfying the fourth-order recurrence
$a(n{+}4) = 2\,a(n{+}3) + 4\,a(n{+}2) - a(n)$ of the prism graph's own matching
sequence A102080 (same recurrence signature, different initial values). For
every $n \ge 3$, the number of matchings of the $n$-prism graph
$C_n \square P_2$ with one spoke (rung) edge deleted equals
$\mathrm{aspoke}(n-3)$, and with one rim (perimeter) edge deleted equals
$\mathrm{arim}(n-3)$; matchings are counted in the Hosoya sense, the empty
matching included. Both theorems are formalised in Lean 4 (mathlib
v4.34.0) and verified by the kernel with no \texttt{sorry}; \texttt{\#print
axioms} reports only \{propext, \texttt{Quot.sound},
\texttt{Classical.choice}\}. The pipeline's exclusivity review found no
occupying record (reachable channels): the two sequences are unregistered and
the graph-side bridge is graded first formalisation (Lean three-layer zero
same-form); the value tier is new sequence, never new mathematics and never a
new recurrence. This note is a priority claim of record, not a full paper.
\end{abstract}

\section{Introduction}
This note records two formalised bridge theorems produced by the AI4Math
automated pipeline, in which LLM workers generate statements and proof scripts
and the Lean~4 kernel is the sole adjudicator at every gate: the statement was
frozen by hash before proving, and a separate audit re-computed the symmetric
difference, scanned the whole file for \texttt{sorry}/\texttt{admit} including
comments, and re-printed the axiom dependencies \cite{lean4}. The object is
graph-side and combinatorial: the matching count (Hosoya index, edge-subset
formulation, empty matching included) of the $n$-prism with a single edge
deleted, for the two edge orbits of the prism (one spoke, one rim). The proof
is a multi-stage assembly (split identity, structural decomposition,
transfer/cell lemmas, strong induction on the sequences) built on four
library modules that are themselves kernel-checked and sorry-free. The note is
a priority claim of record, deposited for timestamping; the corresponding
narrative paper is prepared separately and will reference this deposit.
Verification strength is stated in Section~\ref{sec:verif}. No literature
survey is included---that expectation belongs to arXiv moderation and does not
apply to a Zenodo deposit.

\section{Statement}\label{sec:statement}

We fix notation. The prism graph on $n$ vertices per layer is formalised as
$\mathtt{prism}\ n = \mathtt{cycleGraph}\ n\ \square\ \mathtt{pathGraph}\ 2$
on the vertex type $\texttt{Fin}\ n \times \texttt{Fin}\ 2$. The two defect
slices delete one edge of column 0 with both endpoints kept:
the spoke (rung) $((0,0),(0,1))$ gives $\mathtt{prismDelSpoke}\ n$, and the
rim (perimeter) edge $((0,0),(1,0))$ gives $\mathtt{prismDelRim}\ n$.
\texttt{hosoya} $G$ counts the matchings of $G$ as the number of edge subsets
whose elements pairwise have disjoint endpoint sets, the empty matching
included (Hosoya's definition). The two sequences are defined by the same
fourth-order recurrence with the stated initial values; index $0$ corresponds
to graph-side $n = 3$. The two theorems formalised and proved here are:

% LEAN: tasks/20260929-l14c-graphside-pipeline :: hosoya_prism_del_spoke  grade=完全证明
\begin{theorem}[spoke-side bridge]\label{thm:spoke}
For every $n \in \mathbb{N}$ with $3 \le n$,
\[
(\mathtt{hosoya}\,(\mathtt{prismDelSpoke}\ n) : \mathbb{Z}) = \mathtt{aspoke}\,(n-3).
\]
\end{theorem}

% LEAN: tasks/20260929-l14c-graphside-pipeline :: hosoya_prism_del_rim  grade=完全证明
\begin{theorem}[rim-side bridge]\label{thm:rim}
For every $n \in \mathbb{N}$ with $3 \le n$,
\[
(\mathtt{hosoya}\,(\mathtt{prismDelRim}\ n) : \mathbb{Z}) = \mathtt{arim}\,(n-3).
\]
\end{theorem}

\noindent The subscript $n-3$ is natural-number truncated subtraction, exact
on the hypothesis domain $3 \le n$; the frozen statement transcribes the
frozen natural-language wording verbatim and deliberately keeps this form.

\begin{lstlisting}[caption={Frozen statement snapshot (Lean 4, verbatim lines
36--98 of \texttt{formalized/statement.lean}, sha256 \texttt{d186ff16}), as an
ordered subsequence with exactly three lines excluded, declared immediately
below.},label={lst:stmt}]
-- ===== 基座代码（逐字，禁改） =====

instance (n : ℕ) : DecidableRel (pathGraph n).Adj := fun _ _ =>
  decidable_of_iff _ pathGraph_adj.symm

instance {α β : Type*} [DecidableEq α] [DecidableEq β] (G : SimpleGraph α) (H : SimpleGraph β)
    [DecidableRel G.Adj] [DecidableRel H.Adj] : DecidableRel (G □ H).Adj := fun _ _ =>
  decidable_of_iff _ boxProd_adj.symm

def edgeSupp {V : Type*} [DecidableEq V] [Fintype V] (e : Sym2 V) : Finset V :=
  (univ : Finset V).filter fun v => v ∈ e

def hosoya (G : SimpleGraph V) [Fintype V] [DecidableEq V] [DecidableRel G.Adj] : ℕ :=
  ((G.edgeFinset.powerset).filter fun s =>
    (∀ e ∈ s, ∀ f ∈ s, e ≠ f → Disjoint (edgeSupp e) (edgeSupp f))).card

def prism (n : ℕ) : SimpleGraph (Fin n × Fin 2) := cycleGraph n □ pathGraph 2

instance (n : ℕ) : DecidableRel (prism n).Adj :=
  show DecidableRel (cycleGraph n □ pathGraph 2).Adj from inferInstance

def prismDelSpoke (n : ℕ) (hn : 3 ≤ n) : SimpleGraph (Fin n × Fin 2) :=
  (prism n).deleteEdges {Sym2.mk ((⟨0, by omega⟩, (0 : Fin 2))) ((⟨0, by omega⟩, (1 : Fin 2)))}

instance (n : ℕ) (hn : 3 ≤ n) : DecidableRel (prismDelSpoke n hn).Adj :=
  show DecidableRel ((prism n).deleteEdges
    {Sym2.mk ((⟨0, by omega⟩, (0 : Fin 2))) ((⟨0, by omega⟩, (1 : Fin 2)))}).Adj from inferInstance

def prismDelRim (n : ℕ) (hn : 3 ≤ n) : SimpleGraph (Fin n × Fin 2) :=
  (prism n).deleteEdges {Sym2.mk ((⟨0, by omega⟩, ⟨0, by omega⟩)) ((⟨1, by omega⟩, ⟨0, by omega⟩))}

instance (n : ℕ) (hn : 3 ≤ n) : DecidableRel (prismDelRim n hn).Adj :=
  show DecidableRel ((prism n).deleteEdges
    {Sym2.mk ((⟨0, by omega⟩, ⟨0, by omega⟩)) ((⟨1, by omega⟩, ⟨0, by omega⟩))}).Adj from inferInstance

def aspoke : ℕ → ℤ
  | 0 => 25
  | 1 => 86
  | 2 => 271
  | 3 => 876
  | n + 4 => 2 * aspoke (n + 3) + 4 * aspoke (n + 2) - aspoke n

def arim : ℕ → ℤ
  | 0 => 26
  | 1 => 86
  | 2 => 274
  | 3 => 883
  | n + 4 => 2 * arim (n + 3) + 4 * arim (n + 2) - arim n


/-- **主定理一（spoke 侧桥定理）**：对一切 `n ≥ 3`，删一条 spoke 的棱柱的匹配数等于序列
`aspoke (n - 3)`。其中 `aspoke` 为已冻结序列（初值 25, 86, 271, 876，四阶递推
`a(n+4) = 2·a(n+3) + 4·a(n+2) − a(n)`，下标 0 对应图侧 n = 3）。 -/
theorem hosoya_prism_del_spoke : ∀ (n : ℕ) (hn : 3 ≤ n),
    (hosoya (prismDelSpoke n hn) : ℤ) = aspoke (n - 3) := by

/-- **主定理二（rim 侧桥定理）**：对一切 `n ≥ 3`，删一条 rim 边的棱柱的匹配数等于序列
`arim (n - 3)`。其中 `arim` 为已冻结序列（初值 26, 86, 274, 883，同四阶递推）。 -/
theorem hosoya_prism_del_rim : ∀ (n : ℕ) (hn : 3 ≤ n),
    (hosoya (prismDelRim n hn) : ℤ) = arim (n - 3) := by
\end{lstlisting}

\noindent This is lines 36--98 of the frozen statement file with exactly three
lines excluded, all declared here: the two proof-body placeholder lines
(line 92 and line 98, each the single line \texttt{sorry}, the formalisation
department's placeholders for the two bridge theorems, which this note's
Section~\ref{sec:proof} replaces with the real kernel-checked proofs), and the
one-line separator comment at line 85, which is excluded only because the
paper-lint listing gate rejects the token it contains inside listings and the
line is otherwise a vacuous section marker. Everything else---the header of
the appended block, both doc-comments, all base definitions, both instances
and both theorem statements---is verbatim and in order.
\noindent \texttt{inspect-tex.py}, the package-local checker copied under
\texttt{audit/}, checks this mechanically: every remaining line
occurs in the source in order, and the only lines missing from the block are
those three. The real, sorry-free proofs follow in full.

\section{Proof}\label{sec:proof}

The proof file imports four kernel-checked library modules (the split identity
of the pipeline's Phase 0, the Phase 1 structural lemmas, the ladder-side and
prism-side recurrences) and assembles the two bridges in three lanes. Lane A
(split): a general edge-deletion identity cuts the full prism count along the
deleted edge into the two defect factors plus an endpoint-deleted factor.
Lane B (structure): column-wise decompositions express the two defect counts
through the intact-prism sequence and an auxiliary sequence, packaged as
cell lemmas closed by \texttt{decide} over an eight-step library recurrence
chain. Lane C (induction): the strong-induction closers
\texttt{aspoke\_strong}/\texttt{arim\_strong} derive the sequence values from
the recurrence, and the two frozen theorem heads are discharged from the
$\mathbb{Z}$-cast bridges by \texttt{omega}. The complete file (1857 lines)
follows verbatim; it also ships in this deposit under \texttt{proofs/}.

\begin{lstlisting}[basicstyle=\ttfamily\footnotesize,caption={Complete proof
file of Theorems~\ref{thm:spoke} and~\ref{thm:rim}, verbatim lines 1--1857
(whole file) of \texttt{attempts/phase2/r3/bridge.lean} (sha256
\texttt{399ee81f}). Mixed line endings on the working disk (the file is CRLF
throughout) are spliced LF-normalised per the chorded/ladder/hosoya precedent;
the byte-exact original ships in \texttt{proofs/}.},label={lst:proof}]
import Mathlib
import Library.HosoyaPrismSplit
import Library.GraphsidePhase1
import Library.LadderHosoya
import Library.PrismHosoya

/-!
# Phase 2 r3 · 双桥定理 + 棱柱 R4 收口（bridge.lean，2026-10-01）

provenance：
- 任务契约 `attempts/phase2/r3/prompt.md`（Phase 2 第 4 采样位，最后一个）。
- 基座 def（aspoke/arim/prismDelRim + instance）逐字拷自冻结快照
  `formalized/statement.lean` L64–L83（宪条 2：只许拷贝不许改）。
- 数学骨架（契约 T4a–T4f）：p(n) = a(n)+2c(n)+a(n-2) 四格分解；
  d_rim(n) = a(n)+c(n)（列旋转共轭）；d_spoke(n) = p(n)-a(n-1)（Phase 0 分割
  + 同构）；归纳桥走强归纳 + aspoke/arim 自身 R4（ℤ）。库件 G0 全部经 kernel
  预核验（本文件编译即复验）。
- 战术承 r1/r2 模板：rfl 桥（Fin.val_mk / pair 投影）、omega 断原子先放桥、
  membership 经 `eset_mem_mk`/`lad_adj_iff`/`disjoint_esupp_iff` 线性化。
- n 一律写成 m+3（Fin (m+3) 上字面量 1 与算术实例可用；Fin n 一般形无
  OfNat 实例，r1/r2 同款移位口径）。
-/

open SimpleGraph Finset

-- ===== 基座 def（逐字拷自冻结快照 formalized/statement.lean，勿改） =====

def prismDelRim (n : ℕ) (hn : 3 ≤ n) : SimpleGraph (Fin n × Fin 2) :=
  (prism n).deleteEdges {Sym2.mk ((⟨0, by omega⟩, ⟨0, by omega⟩)) ((⟨1, by omega⟩, ⟨0, by omega⟩))}

instance (n : ℕ) (hn : 3 ≤ n) : DecidableRel (prismDelRim n hn).Adj :=
  show DecidableRel ((prism n).deleteEdges
    {Sym2.mk ((⟨0, by omega⟩, ⟨0, by omega⟩)) ((⟨1, by omega⟩, ⟨0, by omega⟩))}).Adj from inferInstance

def aspoke : ℕ → ℤ
  | 0 => 25
  | 1 => 86
  | 2 => 271
  | 3 => 876
  | n + 4 => 2 * aspoke (n + 3) + 4 * aspoke (n + 2) - aspoke n

def arim : ℕ → ℤ
  | 0 => 26
  | 1 => 86
  | 2 => 274
  | 3 => 883
  | n + 4 => 2 * arim (n + 3) + 4 * arim (n + 2) - arim n

namespace PrismR3

/-! ## §0 小工具（Fin 2 行算术 / Fin (m+3) 模加一 / wrap 边） -/

-- Fin 2 行算术（全枚举 decide）
theorem one_sub_val : ∀ r : Fin 2, ((1 - r : Fin 2) : ℕ) = 1 - (r : ℕ) := by decide
theorem one_sub_ne : ∀ r s : Fin 2, r ≠ s → (1 - r : Fin 2) ≠ 1 - s := by decide
theorem one_sub_eq_one : ∀ r : Fin 2, ((1 - r : Fin 2) : ℕ) = 1 ↔ r = 0 := by decide
theorem one_sub_eq_zero : ∀ r : Fin 2, ((1 - r : Fin 2) : ℕ) = 0 ↔ r = 1 := by decide

-- Fin (m+3) 加一（含环绕）的 val 字符化
theorem fin_add_one_val (m : ℕ) (u v : Fin (m+3)) :
    v = u + 1 ↔ (v : ℕ) = (u:ℕ) + 1 ∨ ((u:ℕ) = m + 2 ∧ (v:ℕ) = 0) := by
  constructor
  · intro h
    have hc := congrArg Fin.val h
    rw [Fin.val_add, Fin.val_one] at hc
    rcases Nat.lt_or_ge ((u:ℕ) + 1) (m + 3) with hlt | hge
    · rw [Nat.mod_eq_of_lt hlt] at hc
      exact Or.inl hc
    · have hu : (u:ℕ) = m + 2 := by omega
      have hmod : ((u:ℕ) + 1) % (m + 3) = 0 := by
        rw [show (u:ℕ) + 1 = m + 3 from by omega]
        exact Nat.mod_self _
      rw [hmod] at hc
      exact Or.inr ⟨hu, hc⟩
  · rintro (h | ⟨hu, hv⟩)
    · apply Fin.val_injective
      rw [Fin.val_add, Fin.val_one]
      rw [show ((u:ℕ) + 1) % (m + 3) = (u:ℕ) + 1 from Nat.mod_eq_of_lt (by omega)]
      exact h
    · apply Fin.val_injective
      rw [Fin.val_add, Fin.val_one]
      rw [show ((u:ℕ) + 1) % (m + 3) = 0 from by
        rw [show (u:ℕ) + 1 = m + 3 from by omega]; exact Nat.mod_self _]
      exact hv

/-- 环图 vs 路图二分：cycleGraph (m+3) 的邻接 = 路邻接 ∪ 环绕边（val 口径）。 -/
theorem cyc_path (m : ℕ) (u v : Fin (m+3)) :
    (cycleGraph (m+3)).Adj u v ↔ (pathGraph (m+3)).Adj u v ∨
      ((u:ℕ) = m + 2 ∧ (v:ℕ) = 0) ∨ ((u:ℕ) = 0 ∧ (v:ℕ) = m + 2) := by
  rw [cycleGraph_adj_iff_succ, pathGraph_adj, fin_add_one_val m u v, fin_add_one_val m v u]
  constructor
  · rintro ((h | ⟨h1, h2⟩) | (h | ⟨h1, h2⟩))
    · exact Or.inl (Or.inl (by omega))
    · exact Or.inr (Or.inl ⟨h1, h2⟩)
    · exact Or.inl (Or.inr (by omega))
    · exact Or.inr (Or.inr ⟨h2, h1⟩)
  · rintro ((h | h) | ⟨h1, h2⟩ | ⟨h1, h2⟩)
    · exact Or.inl (Or.inl (by omega))
    · exact Or.inr (Or.inl (by omega))
    · exact Or.inl (Or.inr ⟨h1, h2⟩)
    · exact Or.inr (Or.inr ⟨h2, h1⟩)

theorem cyc_of_path (m : ℕ) (u v : Fin (m+3)) (h : (pathGraph (m+3)).Adj u v) :
    (cycleGraph (m+3)).Adj u v := (cyc_path m u v).mpr (Or.inl h)

theorem cyc_wrap_0 (m : ℕ) :
    (cycleGraph (m+3)).Adj (⟨m+2, by omega⟩ : Fin (m+3)) (⟨0, by omega⟩ : Fin (m+3)) := by
  refine (cyc_path m _ _).mpr (Or.inr (Or.inl ⟨?_, ?_⟩))
  · exact rfl
  · exact rfl

theorem cyc_wrap_1 (m : ℕ) :
    (cycleGraph (m+3)).Adj (⟨0, by omega⟩ : Fin (m+3)) (⟨m+2, by omega⟩ : Fin (m+3)) := by
  refine (cyc_path m _ _).mpr (Or.inr (Or.inr ⟨?_, ?_⟩))
  · exact rfl
  · exact rfl

theorem top_val (m : ℕ) : ((⟨m+2, by omega⟩ : Fin (m+3)) : ℕ) = m + 2 := rfl
theorem zero_val (m : ℕ) : ((⟨0, by omega⟩ : Fin (m+3)) : ℕ) = 0 := rfl

/-- 环绕 rim 边：第 m+2 列 ~ 第 0 列，行 j。 -/
def wE (m : ℕ) (j : Fin 2) : Sym2 (Fin (m+3) × Fin 2) :=
  s((⟨m+2, by omega⟩, j), (⟨0, by omega⟩, j))

theorem f2_01 : (0 : Fin 2) ≠ (1 : Fin 2) := by decide

theorem wE_nondisj (m : ℕ) (j : Fin 2) :
    ¬ Disjoint (edgeSupp (wE m j)) (edgeSupp (wE m j)) := by
  intro h
  have h4 := (disjoint_esupp_iff ((⟨m+2, by omega⟩, j) : Fin (m+3) × Fin 2)
    ((⟨0, by omega⟩, j) : Fin (m+3) × Fin 2)
    ((⟨m+2, by omega⟩, j) : Fin (m+3) × Fin 2)
    ((⟨0, by omega⟩, j) : Fin (m+3) × Fin 2)).mp h
  exact h4.1 rfl

theorem wE_ne (m : ℕ) : wE m 0 ≠ wE m 1 := by
  intro h
  simp only [wE] at h
  rw [Sym2.eq_iff] at h
  rcases h with ⟨h1, h2⟩ | ⟨h1, h2⟩
  · exact absurd (congrArg Prod.snd h1) f2_01
  · have hf : (⟨m+2, by omega⟩ : Fin (m+3)) = (⟨0, by omega⟩ : Fin (m+3)) :=
      congrArg Prod.fst h1
    have hv : ((⟨m+2, by omega⟩ : Fin (m+3)) : ℕ) = ((⟨0, by omega⟩ : Fin (m+3)) : ℕ) :=
      congrArg Fin.val hf
    rw [top_val, zero_val] at hv
    omega

theorem wE_not_mem_Eset (m : ℕ) (j : Fin 2) : wE m j ∉ Eset (m+3) (m+3) := by
  intro h
  have hadj := (eset_mem_mk ((⟨m+2, by omega⟩, j) : Fin (m+3) × Fin 2)
    ((⟨0, by omega⟩, j) : Fin (m+3) × Fin 2)).mp h
  have h1 := hadj.1
  rw [lad_adj_iff] at h1
  rcases h1 with ⟨hc, hr⟩ | ⟨hr, hc⟩
  · have hA : (((⟨m+2, by omega⟩, j) : Fin (m+3) × Fin 2).1 : ℕ) = m + 2 := rfl
    have hB : (((⟨0, by omega⟩, j) : Fin (m+3) × Fin 2).1 : ℕ) = 0 := rfl
    have hv := congrArg Fin.val hc
    rw [hA, hB] at hv
    omega
  · rcases hc with hc | hc
    · have hA : (((⟨m+2, by omega⟩, j) : Fin (m+3) × Fin 2).1 : ℕ) = m + 2 := rfl
      have hB : (((⟨0, by omega⟩, j) : Fin (m+3) × Fin 2).1 : ℕ) = 0 := rfl
      rw [hA, hB] at hc
      omega
    · have hA : (((⟨m+2, by omega⟩, j) : Fin (m+3) × Fin 2).1 : ℕ) = m + 2 := rfl
      have hB : (((⟨0, by omega⟩, j) : Fin (m+3) × Fin 2).1 : ℕ) = 0 := rfl
      rw [hA, hB] at hc
      omega

/-- Eset (m+3) (m+3) 的成员资格 = 梯子邻接（列界自动成立）。 -/
theorem esetNN (m : ℕ) (a b : Fin (m+3) × Fin 2) :
    s(a, b) ∈ Eset (m+3) (m+3) ↔ (pathGraph (m+3) □ pathGraph 2).Adj a b := by
  rw [eset_mem_mk]
  exact ⟨fun h => h.1, fun h => ⟨h, a.1.isLt, b.1.isLt⟩⟩

/-- Cst (m+3) (m+3) 的成员资格线性化（r2 §1 口径）。 -/
theorem cstNN (m : ℕ) (a b : Fin (m+3) × Fin 2) :
    s(a, b) ∈ PrismR2.Cst (m+3) (m+3) ↔ (pathGraph (m+3) □ pathGraph 2).Adj a b ∧
      ((a.2 : ℕ) = 1 ∨ 1 ≤ (a.1 : ℕ) ∧ (a.1 : ℕ) + 2 ≤ m + 3) ∧
      ((b.2 : ℕ) = 1 ∨ 1 ≤ (b.1 : ℕ) ∧ (b.1 : ℕ) + 2 ≤ m + 3) := by
  simp only [PrismR2.Cst, Finset.mem_filter, PrismR2.filt2, esetNN]

/-- 与 wE 的 disjoint 条件（列/行线性口径）。 -/
theorem disj_wE_iff (m : ℕ) (j : Fin 2) (a b : Fin (m+3) × Fin 2) :
    Disjoint (edgeSupp (s(a, b))) (edgeSupp (wE m j)) ↔
      ((a.1 : ℕ) ≠ m + 2 ∨ a.2 ≠ j) ∧ ((a.1 : ℕ) ≠ 0 ∨ a.2 ≠ j) ∧
      ((b.1 : ℕ) ≠ m + 2 ∨ b.2 ≠ j) ∧ ((b.1 : ℕ) ≠ 0 ∨ b.2 ≠ j) := by
  have h := disjoint_esupp_iff a b (⟨m+2, by omega⟩, j) (⟨0, by omega⟩, j)
  constructor
  · intro hd
    obtain ⟨h1, h2, h3, h4⟩ := h.mp hd
    refine ⟨?_, ?_, ?_, ?_⟩
    · simpa only [top_val] using PrismR2.ne_to_col _ _ _ h1
    · simpa only [zero_val] using PrismR2.ne_to_col _ _ _ h2
    · simpa only [top_val] using PrismR2.ne_to_col _ _ _ h3
    · simpa only [zero_val] using PrismR2.ne_to_col _ _ _ h4
  · rintro ⟨h1, h2, h3, h4⟩
    refine h.mpr ⟨?_, ?_, ?_, ?_⟩
    · exact PrismR2.col_to_ne _ _ _ (by simpa only [top_val] using h1)
    · exact PrismR2.col_to_ne _ _ _ (by simpa only [zero_val] using h2)
    · exact PrismR2.col_to_ne _ _ _ (by simpa only [top_val] using h3)
    · exact PrismR2.col_to_ne _ _ _ (by simpa only [zero_val] using h4)

/-! ## §0b 三个结构映射：swapv（行翻转对合）/ rotv（列旋转）/ shEmb（列平移） -/

/-- 行翻转对合。 -/
def swapv {N : ℕ} (v : Fin N × Fin 2) : Fin N × Fin 2 := (v.1, 1 - v.2)

theorem swapv_inv {N : ℕ} (v : Fin N × Fin 2) : swapv (swapv v) = v := by
  have hkey : (1 - (1 - v.2) : Fin 2) = v.2 := by
    rcases fin2_ex v.2 with h | h <;> rw [h] <;> rfl
  show (v.1, 1 - (1 - v.2)) = v
  rw [hkey]

theorem swapv_inj {N : ℕ} : Function.Injective (swapv : Fin N × Fin 2 → Fin N × Fin 2) := by
  intro v v' h
  have h2 : swapv (swapv v) = swapv (swapv v') := by rw [h]
  rw [swapv_inv, swapv_inv] at h2
  exact h2

theorem swapv_fst {N : ℕ} (v : Fin N × Fin 2) : (swapv v).1 = v.1 := rfl
theorem swapv_fst_val {N : ℕ} (v : Fin N × Fin 2) : ((swapv v).1 : ℕ) = (v.1 : ℕ) := rfl
theorem swapv_snd_val {N : ℕ} (v : Fin N × Fin 2) : ((swapv v).2 : ℕ) = 1 - (v.2 : ℕ) :=
  one_sub_val v.2

/-- 列旋转（+1 mod m+3），行不变。 -/
def rotv (m : ℕ) (v : Fin (m+3) × Fin 2) : Fin (m+3) × Fin 2 := (v.1 + 1, v.2)

theorem rotv_val_lt (m : ℕ) (v : Fin (m+3) × Fin 2) (h : (v.1 : ℕ) + 1 < m + 3) :
    ((rotv m v).1 : ℕ) = (v.1 : ℕ) + 1 := by
  show (((v.1 + 1 : Fin (m+3))) : ℕ) = (v.1 : ℕ) + 1
  rw [Fin.val_add, Fin.val_one]
  exact Nat.mod_eq_of_lt h

theorem rotv_val_top (m : ℕ) (v : Fin (m+3) × Fin 2) (h : (v.1 : ℕ) = m + 2) :
    ((rotv m v).1 : ℕ) = 0 := by
  show (((v.1 + 1 : Fin (m+3))) : ℕ) = 0
  rw [Fin.val_add, Fin.val_one]
  have hmod : ((v.1 : ℕ) + 1) % (m + 3) = 0 := by
    rw [show (v.1 : ℕ) + 1 = m + 3 from by omega]
    exact Nat.mod_self _
  rw [hmod]

theorem rotv_row (m : ℕ) (v : Fin (m+3) × Fin 2) : (rotv m v).2 = v.2 := rfl

theorem rotv_inj (m : ℕ) : Function.Injective (rotv m) := by
  intro v v' h
  have hrow : (rotv m v).2 = (rotv m v').2 := congrArg (fun p : Fin (m+3) × Fin 2 => p.2) h
  have hcol0 : (rotv m v).1 = (rotv m v').1 := congrArg (fun p : Fin (m+3) × Fin 2 => p.1) h
  have ha : (rotv m v).1 = v.1 + 1 := rfl
  have hb : (rotv m v').1 = v'.1 + 1 := rfl
  rw [ha, hb] at hcol0
  have hval := congrArg Fin.val hcol0
  rw [Fin.val_add, Fin.val_one, Fin.val_add, Fin.val_one] at hval
  have hv : (v.1 : ℕ) < m + 3 := v.1.isLt
  have hv' : (v'.1 : ℕ) < m + 3 := v'.1.isLt
  have hcc : (v.1 : ℕ) = (v'.1 : ℕ) := by
    rcases Nat.lt_or_ge ((v.1 : ℕ) + 1) (m + 3) with h1 | h1
    · rcases Nat.lt_or_ge ((v'.1 : ℕ) + 1) (m + 3) with h2 | h2
      · rw [Nat.mod_eq_of_lt h1, Nat.mod_eq_of_lt h2] at hval
        omega
      · rw [Nat.mod_eq_of_lt h1] at hval
        have hmod : ((v'.1 : ℕ) + 1) % (m + 3) = 0 := by
          rw [show (v'.1 : ℕ) + 1 = m + 3 from by omega]
          exact Nat.mod_self _
        rw [hmod] at hval
        omega
    · rcases Nat.lt_or_ge ((v'.1 : ℕ) + 1) (m + 3) with h2 | h2
      · rw [Nat.mod_eq_of_lt h2] at hval
        have hmod : ((v.1 : ℕ) + 1) % (m + 3) = 0 := by
          rw [show (v.1 : ℕ) + 1 = m + 3 from by omega]
          exact Nat.mod_self _
        rw [hmod] at hval
        omega
      · have hmod1 : ((v.1 : ℕ) + 1) % (m + 3) = 0 := by
          rw [show (v.1 : ℕ) + 1 = m + 3 from by omega]
          exact Nat.mod_self _
        have hmod2 : ((v'.1 : ℕ) + 1) % (m + 3) = 0 := by
          rw [show (v'.1 : ℕ) + 1 = m + 3 from by omega]
          exact Nat.mod_self _
        rw [hmod1, hmod2] at hval
        omega
  exact Prod.ext_iff.mpr ⟨Fin.val_injective hcc, by exact hrow⟩

/-- 列平移嵌入（col ↦ col+1）：Fin (m+1) × Fin 2 → Fin (m+3) × Fin 2。 -/
def shEmb (m : ℕ) (v : Fin (m+1) × Fin 2) : Fin (m+3) × Fin 2 :=
  (⟨(v.1 : ℕ) + 1, by omega⟩, v.2)

theorem shEmb_inj (m : ℕ) : Function.Injective (shEmb m) := by
  intro v v' h
  have hrow : (shEmb m v).2 = (shEmb m v').2 := by rw [h]
  have hcol : ((shEmb m v).1 : ℕ) = ((shEmb m v').1 : ℕ) := by rw [h]
  have hA : ((shEmb m v).1 : ℕ) = (v.1 : ℕ) + 1 := rfl
  have hB : ((shEmb m v').1 : ℕ) = (v'.1 : ℕ) + 1 := rfl
  rw [hA, hB] at hcol
  have hv : (v.1 : ℕ) < m + 1 := v.1.isLt
  have hv' : (v'.1 : ℕ) < m + 1 := v'.1.isLt
  exact Prod.ext_iff.mpr ⟨Fin.val_injective (by omega), by exact hrow⟩

theorem shEmb_val (m : ℕ) (v : Fin (m+1) × Fin 2) :
    ((shEmb m v).1 : ℕ) = (v.1 : ℕ) + 1 := rfl
theorem shEmb_row (m : ℕ) (v : Fin (m+1) × Fin 2) : (shEmb m v).2 = v.2 := rfl

/-- emb 单射（r1 transp 证明体内嵌的核验件，抽出复用）。 -/
theorem emb_inj {N k : ℕ} (hk : k ≤ N) : Function.Injective (emb hk) := by
  intro v v' h
  have h1 : (⟨(v.1 : ℕ), by omega⟩ : Fin N) = (⟨(v'.1 : ℕ), by omega⟩ : Fin N) :=
    congrArg Prod.fst h
  have hv : v.1 = v'.1 := by
    apply Fin.val_injective
    show (v.1 : ℕ) = (v'.1 : ℕ)
    exact congrArg (fun i : Fin N => (i : ℕ)) h1
  have hs : (emb hk v).2 = (emb hk v').2 := congrArg (fun p : Fin N × Fin 2 => p.2) h
  exact Prod.ext_iff.mpr ⟨hv, hs⟩

/-! ## §1 T4a：棱柱边集两分类（梯子边集 + 双环绕 rim 边） -/

theorem s_eq_wE_iff (m : ℕ) (j : Fin 2) (a b : Fin (m+3) × Fin 2) :
    s(a, b) = wE m j ↔
      (a = (⟨m+2, by omega⟩, j) ∧ b = (⟨0, by omega⟩, j)) ∨
      (a = (⟨0, by omega⟩, j) ∧ b = (⟨m+2, by omega⟩, j)) := by
  simp only [wE, Sym2.eq_iff]

theorem s_eq_wE_of (m : ℕ) (j : Fin 2) (a b : Fin (m+3) × Fin 2)
    (ha : (a.1 : ℕ) = m + 2) (hb : (b.1 : ℕ) = 0) (hja : a.2 = j) (hjb : b.2 = j) :
    s(a, b) = wE m j := by
  rw [s_eq_wE_iff]
  exact Or.inl ⟨Prod.ext_iff.mpr ⟨Fin.val_injective ha, hja⟩,
    Prod.ext_iff.mpr ⟨Fin.val_injective hb, hjb⟩⟩

theorem s_eq_wE_of' (m : ℕ) (j : Fin 2) (a b : Fin (m+3) × Fin 2)
    (ha : (a.1 : ℕ) = 0) (hb : (b.1 : ℕ) = m + 2) (hja : a.2 = j) (hjb : b.2 = j) :
    s(a, b) = wE m j := by
  rw [s_eq_wE_iff]
  exact Or.inr ⟨Prod.ext_iff.mpr ⟨Fin.val_injective ha, hja⟩,
    Prod.ext_iff.mpr ⟨Fin.val_injective hb, hjb⟩⟩

theorem wE_adj_prism (m : ℕ) (j : Fin 2) :
    (prism (m+3)).Adj (⟨m+2, by omega⟩, j) (⟨0, by omega⟩, j) := by
  show (cycleGraph (m+3) □ pathGraph 2).Adj (⟨m+2, by omega⟩, j) (⟨0, by omega⟩, j)
  exact boxProd_adj.mpr (Or.inl ⟨cyc_wrap_0 m, rfl⟩)

theorem wE_adj_prism' (m : ℕ) (j : Fin 2) :
    (prism (m+3)).Adj (⟨0, by omega⟩, j) (⟨m+2, by omega⟩, j) := by
  show (cycleGraph (m+3) □ pathGraph 2).Adj (⟨0, by omega⟩, j) (⟨m+2, by omega⟩, j)
  exact boxProd_adj.mpr (Or.inl ⟨cyc_wrap_1 m, rfl⟩)

theorem wE_mem_prism (m : ℕ) (j : Fin 2) : wE m j ∈ (prism (m+3)).edgeFinset := by
  simp only [wE]
  rw [SimpleGraph.mem_edgeFinset, SimpleGraph.mem_edgeSet]
  exact wE_adj_prism m j

/-- **T4a**：棱柱边集 = Eset（梯子）∪ 双环绕。 -/
theorem prism_edges (m : ℕ) :
    (prism (m+3)).edgeFinset = Eset (m+3) (m+3) ∪ {wE m 0, wE m 1} := by
  ext e
  induction e using Sym2.ind with
  | _ a b =>
  simp only [Finset.mem_union, Finset.mem_insert, Finset.mem_singleton, esetNN,
    SimpleGraph.mem_edgeFinset, SimpleGraph.mem_edgeSet]
  have hb : (cycleGraph (m+3) □ pathGraph 2).Adj a b ↔
      ((cycleGraph (m+3)).Adj a.1 b.1 ∧ a.2 = b.2) ∨
      ((pathGraph 2).Adj a.2 b.2 ∧ a.1 = b.1) := boxProd_adj
  constructor
  · intro h
    have h' : (cycleGraph (m+3) □ pathGraph 2).Adj a b := h
    rcases hb.mp h' with ⟨hcyc, hrow⟩ | ⟨hp2, hcol⟩
    · rcases (cyc_path m a.1 b.1).mp hcyc with hpath | ⟨h1, h2⟩ | ⟨h1, h2⟩
      · exact Or.inl (by rw [lad_adj_iff]; exact Or.inr ⟨hrow, pathGraph_adj.mp hpath⟩)
      · rcases fin2_ex a.2 with hr | hr
        · exact Or.inr (Or.inl (s_eq_wE_of m 0 a b h1 h2 hr (hrow.symm.trans hr)))
        · exact Or.inr (Or.inr (s_eq_wE_of m 1 a b h1 h2 hr (hrow.symm.trans hr)))
      · rcases fin2_ex a.2 with hr | hr
        · exact Or.inr (Or.inl (s_eq_wE_of' m 0 a b h1 h2 hr (hrow.symm.trans hr)))
        · exact Or.inr (Or.inr (s_eq_wE_of' m 1 a b h1 h2 hr (hrow.symm.trans hr)))
    · exact Or.inl (by rw [lad_adj_iff]; exact Or.inl ⟨hcol, (p2adj a.2 b.2).mp hp2⟩)
  · rintro (hlad | hw | hw)
    · rcases (lad_adj_iff.mp hlad) with ⟨hcol, hrow⟩ | ⟨hrow, hc⟩
      · exact boxProd_adj.mpr (Or.inr ⟨(p2adj a.2 b.2).mpr hrow, hcol⟩)
      · exact boxProd_adj.mpr (Or.inl ⟨cyc_of_path m a.1 b.1 (pathGraph_adj.mpr hc), hrow⟩)
    · rw [s_eq_wE_iff] at hw
      rcases hw with ⟨ha, hb2⟩ | ⟨ha, hb2⟩
      · rw [ha, hb2]; exact wE_adj_prism m 0
      · rw [ha, hb2]; exact wE_adj_prism' m 0
    · rw [s_eq_wE_iff] at hw
      rcases hw with ⟨ha, hb2⟩ | ⟨ha, hb2⟩
      · rw [ha, hb2]; exact wE_adj_prism m 1
      · rw [ha, hb2]; exact wE_adj_prism' m 1

theorem hosoya_prism_mc (m : ℕ) :
    hosoya (prism (m+3)) = matchCount (Eset (m+3) (m+3) ∪ {wE m 0, wE m 1}) := by
  rw [hosoya_matchCount, prism_edges]

/-! ## §2 T4b 四格 cell 引理 -/

theorem f2_val0 : ∀ j : Fin 2, (j : ℕ) = 0 → j = 0 := by decide
theorem f2_val1 : ∀ j : Fin 2, (j : ℕ) = 1 → j = 1 := by decide
theorem f2_val_cases : ∀ j : Fin 2, (j : ℕ) = 0 ∨ (j : ℕ) = 1 := by decide
theorem f2_ne01 : ∀ j : Fin 2, (j : ℕ) = 1 → j ≠ 0 := by decide
theorem f2_ne10 : ∀ j : Fin 2, (j : ℕ) = 0 → j ≠ 1 := by decide
theorem one_sub_inj : ∀ r s : Fin 2, (1 - r : Fin 2) = 1 - s → r = s := by decide

theorem lad_adj_swap {N : ℕ} (x y : Fin N × Fin 2) :
    (pathGraph N □ pathGraph 2).Adj (swapv x) (swapv y) ↔
      (pathGraph N □ pathGraph 2).Adj x y := by
  rw [lad_adj_iff, lad_adj_iff]
  constructor
  · rintro (⟨hcol, hr⟩ | ⟨hrow, hc⟩)
    · exact Or.inl ⟨hcol, fun h => hr (congrArg (fun j : Fin 2 => 1 - j) h)⟩
    · exact Or.inr ⟨one_sub_inj x.2 y.2 hrow, hc⟩
  · rintro (⟨hcol, hr⟩ | ⟨hrow, hc⟩)
    · exact Or.inl ⟨hcol, fun h => hr (one_sub_inj x.2 y.2 h)⟩
    · exact Or.inr ⟨congrArg (fun j : Fin 2 => 1 - j) hrow, hc⟩

/-- **cell C2'（T4b·Eset 格）**：Eset 过滤 wE0-不相交格 = Cst。 -/
theorem cell_wE0 (m : ℕ) :
    (Eset (m+3) (m+3)).filter (fun f => Disjoint (edgeSupp f) (edgeSupp (wE m 0))) =
      PrismR2.Cst (m+3) (m+3) := by
  ext e
  induction e using Sym2.ind with
  | _ a b =>
  simp only [Finset.mem_filter, esetNN, cstNN, disj_wE_iff m 0 a b]
  have hais : (a.1 : ℕ) < m + 3 := a.1.isLt
  have hbis : (b.1 : ℕ) < m + 3 := b.1.isLt
  constructor
  · rintro ⟨hadj, h1, h2, h3, h4⟩
    refine ⟨hadj, ?_, ?_⟩
    · rcases f2_val_cases a.2 with hva | hva
      · refine Or.inr ⟨?_, ?_⟩
        · rcases h2 with h | h
          · omega
          · exact absurd (f2_val0 a.2 hva) h
        · rcases h1 with h | h
          · omega
          · exact absurd (f2_val0 a.2 hva) h
      · exact Or.inl hva
    · rcases f2_val_cases b.2 with hvb | hvb
      · refine Or.inr ⟨?_, ?_⟩
        · rcases h4 with h | h
          · omega
          · exact absurd (f2_val0 b.2 hvb) h
        · rcases h3 with h | h
          · omega
          · exact absurd (f2_val0 b.2 hvb) h
      · exact Or.inl hvb
  · rintro ⟨hadj, hqa, hqb⟩
    refine ⟨hadj, ?_, ?_, ?_, ?_⟩
    · rcases f2_val_cases a.2 with hva | hva
      · rcases hqa with hq | hq <;> omega
      · exact Or.inr (f2_ne01 a.2 hva)
    · rcases f2_val_cases a.2 with hva | hva
      · rcases hqa with hq | hq <;> omega
      · exact Or.inr (f2_ne01 a.2 hva)
    · rcases f2_val_cases b.2 with hvb | hvb
      · rcases hqb with hq | hq <;> omega
      · exact Or.inr (f2_ne01 b.2 hvb)
    · rcases f2_val_cases b.2 with hvb | hvb
      · rcases hqb with hq | hq <;> omega
      · exact Or.inr (f2_ne01 b.2 hvb)

theorem f2_flip_ne1 {N : ℕ} (a x : Fin N × Fin 2) (h : (a.2 : ℕ) = 1 - (x.2 : ℕ))
    (hv : (x.2 : ℕ) = 1) : a.2 ≠ (1 : Fin 2) := by
  intro hcon
  have hva : (a.2 : ℕ) = 1 := by rw [hcon]; rfl
  omega

theorem f2_flip_ne1' {N : ℕ} (a x : Fin N × Fin 2) (h : (a.2 : ℕ) = 1 - (x.2 : ℕ))
    (hv : (x.2 : ℕ) = 0) : a.2 ≠ 0 := by
  intro hcon
  have hva : (a.2 : ℕ) = 0 := by rw [hcon]; rfl
  omega

/-- **cell C1（T4b·Eset 格）**：Eset 过滤 wE1-不相交格 = Cst 的行翻转像。 -/
theorem cell_wE1 (m : ℕ) :
    (Eset (m+3) (m+3)).filter (fun f => Disjoint (edgeSupp f) (edgeSupp (wE m 1))) =
      (PrismR2.Cst (m+3) (m+3)).image
        (Sym2.map (swapv : Fin (m+3) × Fin 2 → Fin (m+3) × Fin 2)) := by
  ext e
  induction e using Sym2.ind with
  | _ a b =>
  simp only [Finset.mem_filter, Finset.mem_image, esetNN]
  constructor
  · intro h
    have hadj := h.1
    obtain ⟨h1, h2, h3, h4⟩ := (disj_wE_iff m 1 a b).mp h.2
    refine ⟨s(swapv a, swapv b), ?_, ?_⟩
    · rw [cstNN]
      refine ⟨(lad_adj_swap a b).mpr hadj, ?_, ?_⟩
      · rcases f2_val_cases a.2 with hva | hva
        · exact Or.inl (by have := swapv_snd_val a; omega)
        · refine Or.inr ⟨?_, ?_⟩
          · rcases h2 with h | h
            · have := swapv_fst_val a; omega
            · exact absurd (f2_val1 a.2 hva) h
          · rcases h1 with h | h
            · have := swapv_fst_val a; have := a.1.isLt; omega
            · exact absurd (f2_val1 a.2 hva) h
      · rcases f2_val_cases b.2 with hvb | hvb
        · exact Or.inl (by have := swapv_snd_val b; omega)
        · refine Or.inr ⟨?_, ?_⟩
          · rcases h4 with h | h
            · have := swapv_fst_val b; omega
            · exact absurd (f2_val1 b.2 hvb) h
          · rcases h3 with h | h
            · have := swapv_fst_val b; have := b.1.isLt; omega
            · exact absurd (f2_val1 b.2 hvb) h
    · rw [Sym2.map_mk, swapv_inv, swapv_inv]
  · rintro ⟨e', he', hmap⟩
    induction e' using Sym2.ind with
    | _ x y =>
    rw [Sym2.map_mk, Sym2.eq_iff] at hmap
    obtain ⟨hadj, hqx, hqy⟩ := (cstNN m x y).mp he'
    rcases hmap with ⟨hsa, hsb⟩ | ⟨hsb, hsa⟩
    · have hA : a = swapv x := hsa.symm
      have hB : b = swapv y := hsb.symm
      subst hA
      subst hB
      have ha1 : ((swapv x).1 : ℕ) = (x.1 : ℕ) := swapv_fst_val x
      have ha2 : ((swapv x).2 : ℕ) = 1 - (x.2 : ℕ) := swapv_snd_val x
      have hb1 : ((swapv y).1 : ℕ) = (y.1 : ℕ) := swapv_fst_val y
      have hb2 : ((swapv y).2 : ℕ) = 1 - (y.2 : ℕ) := swapv_snd_val y
      refine ⟨(lad_adj_swap x y).mpr hadj, (disj_wE_iff m 1 (swapv x) (swapv y)).mpr ⟨?_, ?_, ?_, ?_⟩⟩
      · rcases f2_val_cases x.2 with hv | hv
        · left; omega
        · exact Or.inr (f2_flip_ne1 (swapv x) x ha2 hv)
      · rcases f2_val_cases x.2 with hv | hv
        · left; omega
        · exact Or.inr (f2_flip_ne1 (swapv x) x ha2 hv)
      · rcases f2_val_cases y.2 with hv | hv
        · left; omega
        · exact Or.inr (f2_flip_ne1 (swapv y) y hb2 hv)
      · rcases f2_val_cases y.2 with hv | hv
        · left; omega
        · exact Or.inr (f2_flip_ne1 (swapv y) y hb2 hv)
    · have hA : a = swapv y := hsa.symm
      have hB : b = swapv x := hsb.symm
      subst hA
      subst hB
      have ha1 : ((swapv y).1 : ℕ) = (y.1 : ℕ) := swapv_fst_val y
      have ha2 : ((swapv y).2 : ℕ) = 1 - (y.2 : ℕ) := swapv_snd_val y
      have hb1 : ((swapv x).1 : ℕ) = (x.1 : ℕ) := swapv_fst_val x
      have hb2 : ((swapv x).2 : ℕ) = 1 - (x.2 : ℕ) := swapv_snd_val x
      refine ⟨((lad_adj_swap x y).mpr hadj).symm, (disj_wE_iff m 1 (swapv y) (swapv x)).mpr ⟨?_, ?_, ?_, ?_⟩⟩
      · rcases f2_val_cases y.2 with hv | hv
        · left; omega
        · exact Or.inr (f2_flip_ne1 (swapv y) y ha2 hv)
      · rcases f2_val_cases y.2 with hv | hv
        · left; omega
        · exact Or.inr (f2_flip_ne1 (swapv y) y ha2 hv)
      · rcases f2_val_cases x.2 with hv | hv
        · left; omega
        · exact Or.inr (f2_flip_ne1 (swapv x) x hb2 hv)
      · rcases f2_val_cases x.2 with hv | hv
        · left; omega
        · exact Or.inr (f2_flip_ne1 (swapv x) x hb2 hv)

theorem eset_full (N : ℕ) (a b : Fin N × Fin 2) :
    s(a, b) ∈ Eset N N ↔ (pathGraph N □ pathGraph 2).Adj a b := by
  rw [eset_mem_mk]
  exact ⟨fun h => h.1, fun h => ⟨h, a.1.isLt, b.1.isLt⟩⟩

/-- omega 断原子桥（↑v.1 投影原子 omega 不直接连，先抽象成纯 ℕ 变量）。 -/
theorem sub_lt_aux (m c : ℕ) (h1 : 1 ≤ c) (h2 : c ≤ m + 1) : c - 1 < m + 1 := by omega

/-- 列前移一格（col ↦ col-1）的预像函数（1 ≤ col ≤ m+1）。 -/
def vpred (m : ℕ) (v : Fin (m+3) × Fin 2) (h : 1 ≤ (v.1 : ℕ) ∧ (v.1 : ℕ) ≤ m + 1) :
    Fin (m+1) × Fin 2 :=
  (⟨(v.1 : ℕ) - 1, sub_lt_aux m (v.1 : ℕ) h.1 h.2⟩, v.2)

theorem vpred_fst_val (m : ℕ) (v : Fin (m+3) × Fin 2) (h : 1 ≤ (v.1 : ℕ) ∧ (v.1 : ℕ) ≤ m + 1) :
    ((vpred m v h).1 : ℕ) = (v.1 : ℕ) - 1 := rfl
theorem vpred_snd (m : ℕ) (v : Fin (m+3) × Fin 2) (h : 1 ≤ (v.1 : ℕ) ∧ (v.1 : ℕ) ≤ m + 1) :
    (vpred m v h).2 = v.2 := rfl

theorem cst_bound (m : ℕ) (v : Fin (m+3) × Fin 2)
    (hq : (v.2 : ℕ) = 1 ∨ 1 ≤ (v.1 : ℕ) ∧ (v.1 : ℕ) + 2 ≤ m + 3)
    (hc1 : (v.1 : ℕ) ≠ m + 2 ∨ v.2 ≠ (1 : Fin 2))
    (hc2 : (v.1 : ℕ) ≠ 0 ∨ v.2 ≠ (1 : Fin 2)) :
    1 ≤ (v.1 : ℕ) ∧ (v.1 : ℕ) ≤ m + 1 := by
  have his := v.1.isLt
  rcases f2_val_cases v.2 with hv | hv
  · rcases hq with hq' | hq' <;> omega
  · rcases hc1 with h | h
    · rcases hc2 with h' | h'
      · omega
      · exact absurd (f2_val1 v.2 hv) h'
    · exact absurd (f2_val1 v.2 hv) h

theorem lad_adj_pred (m : ℕ) (a b : Fin (m+3) × Fin 2)
    (ha : 1 ≤ (a.1 : ℕ) ∧ (a.1 : ℕ) ≤ m + 1) (hb : 1 ≤ (b.1 : ℕ) ∧ (b.1 : ℕ) ≤ m + 1) :
    (pathGraph (m+1) □ pathGraph 2).Adj (vpred m a ha) (vpred m b hb) ↔
      (pathGraph (m+3) □ pathGraph 2).Adj a b := by
  rw [lad_adj_iff, lad_adj_iff]
  have hxa : ((vpred m a ha).1 : ℕ) = (a.1 : ℕ) - 1 := rfl
  have hxb : ((vpred m b hb).1 : ℕ) = (b.1 : ℕ) - 1 := rfl
  constructor
  · rintro (⟨hc, hr⟩ | ⟨hr, hc⟩)
    · refine Or.inl ⟨?_, hr⟩
      apply Fin.val_injective
      have hcv : (a.1 : ℕ) - 1 = (b.1 : ℕ) - 1 := congrArg Fin.val hc
      omega
    · refine Or.inr ⟨hr, ?_⟩
      rcases hc with hc | hc
      · left; omega
      · right; omega
  · rintro (⟨hc, hr⟩ | ⟨hr, hc⟩)
    · refine Or.inl ⟨?_, hr⟩
      apply Fin.val_injective
      have hcv : (a.1 : ℕ) = (b.1 : ℕ) := congrArg Fin.val hc
      omega
    · refine Or.inr ⟨hr, ?_⟩
      rcases hc with hc | hc
      · left; omega
      · right; omega

theorem shEmb_vpred (m : ℕ) (v : Fin (m+3) × Fin 2) (h : 1 ≤ (v.1 : ℕ) ∧ (v.1 : ℕ) ≤ m + 1) :
    shEmb m (vpred m v h) = v := by
  have hval : ((shEmb m (vpred m v h)).1 : ℕ) = (v.1 : ℕ) := by
    show ((vpred m v h).1 : ℕ) + 1 = (v.1 : ℕ)
    have hx : ((vpred m v h).1 : ℕ) = (v.1 : ℕ) - 1 := rfl
    omega
  exact Prod.ext_iff.mpr ⟨Fin.val_injective hval, rfl⟩

theorem shEmb_row_val (m : ℕ) (v : Fin (m+1) × Fin 2) : ((shEmb m v).2 : ℕ) = (v.2 : ℕ) := rfl

theorem lad_adj_emb (m : ℕ) (x y : Fin (m+1) × Fin 2) :
    (pathGraph (m+3) □ pathGraph 2).Adj (shEmb m x) (shEmb m y) ↔
      (pathGraph (m+1) □ pathGraph 2).Adj x y := by
  rw [lad_adj_iff, lad_adj_iff]
  have hxa : ((shEmb m x).1 : ℕ) = (x.1 : ℕ) + 1 := rfl
  have hxb : ((shEmb m y).1 : ℕ) = (y.1 : ℕ) + 1 := rfl
  have hx : (x.1 : ℕ) < m + 1 := x.1.isLt
  have hy : (y.1 : ℕ) < m + 1 := y.1.isLt
  constructor
  · rintro (⟨hc, hr⟩ | ⟨hr, hc⟩)
    · refine Or.inl ⟨?_, hr⟩
      apply Fin.val_injective
      have hcv : ((shEmb m x).1 : ℕ) = ((shEmb m y).1 : ℕ) := congrArg Fin.val hc
      omega
    · refine Or.inr ⟨hr, ?_⟩
      rcases hc with hc | hc
      · left; omega
      · right; omega
  · rintro (⟨hc, hr⟩ | ⟨hr, hc⟩)
    · refine Or.inl ⟨?_, hr⟩
      apply Fin.val_injective
      have hcv : (x.1 : ℕ) = (y.1 : ℕ) := congrArg Fin.val hc
      omega
    · refine Or.inr ⟨hr, ?_⟩
      rcases hc with hc | hc
      · left; omega
      · right; omega

/-- **cell C3（T4b·Cst 格）**：Cst 过滤 wE1-不相交格 = 列平移梯子像。 -/
theorem cell_C3 (m : ℕ) :
    (PrismR2.Cst (m+3) (m+3)).filter (fun f => Disjoint (edgeSupp f) (edgeSupp (wE m 1))) =
      (Eset (m+1) (m+1)).image (Sym2.map (shEmb m)) := by
  ext e
  induction e using Sym2.ind with
  | _ a b =>
  simp only [Finset.mem_filter, Finset.mem_image, cstNN]
  constructor
  · rintro ⟨⟨hadj, hqa, hqb⟩, hdisj⟩
    obtain ⟨h1, h2, h3, h4⟩ := (disj_wE_iff m 1 a b).mp hdisj
    obtain ⟨hal, har⟩ := cst_bound m a hqa h1 h2
    obtain ⟨hbl, hbr⟩ := cst_bound m b hqb h3 h4
    refine ⟨s(vpred m a ⟨hal, har⟩, vpred m b ⟨hbl, hbr⟩), ?_, ?_⟩
    · rw [eset_full]
      exact (lad_adj_pred m a b ⟨hal, har⟩ ⟨hbl, hbr⟩).mpr hadj
    · rw [Sym2.map_mk, Sym2.eq_iff]
      exact Or.inl ⟨shEmb_vpred m a ⟨hal, har⟩, shEmb_vpred m b ⟨hbl, hbr⟩⟩
  · rintro ⟨e', he', hmap⟩
    induction e' using Sym2.ind with
    | _ x y =>
    rw [Sym2.map_mk, Sym2.eq_iff] at hmap
    have hadj := (eset_full (m+1) x y).mp he'
    rcases hmap with ⟨hsa, hsb⟩ | ⟨hsb, hsa⟩
    · have hA : a = shEmb m x := hsa.symm
      have hB : b = shEmb m y := hsb.symm
      subst hA
      subst hB
      have hxa : ((shEmb m x).1 : ℕ) = (x.1 : ℕ) + 1 := rfl
      have hxb : ((shEmb m y).1 : ℕ) = (y.1 : ℕ) + 1 := rfl
      have hx2 : ((shEmb m x).2 : ℕ) = (x.2 : ℕ) := rfl
      have hy2 : ((shEmb m y).2 : ℕ) = (y.2 : ℕ) := rfl
      have hx : (x.1 : ℕ) < m + 1 := x.1.isLt
      have hy : (y.1 : ℕ) < m + 1 := y.1.isLt
      refine ⟨⟨(lad_adj_emb m x y).mpr hadj, ?_, ?_⟩,
        (disj_wE_iff m 1 (shEmb m x) (shEmb m y)).mpr ⟨?_, ?_, ?_, ?_⟩⟩
      · rcases f2_val_cases x.2 with hv | hv
        · refine Or.inr ⟨?_, ?_⟩
          · omega
          · omega
        · exact Or.inl (by omega)
      · rcases f2_val_cases y.2 with hv | hv
        · refine Or.inr ⟨?_, ?_⟩
          · omega
          · omega
        · exact Or.inl (by omega)
      · rcases f2_val_cases x.2 with hv | hv
        · exact Or.inr (f2_ne10 x.2 hv)
        · left; omega
      · rcases f2_val_cases x.2 with hv | hv
        · exact Or.inr (f2_ne10 x.2 hv)
        · left; omega
      · rcases f2_val_cases y.2 with hv | hv
        · exact Or.inr (f2_ne10 y.2 hv)
        · left; omega
      · rcases f2_val_cases y.2 with hv | hv
        · exact Or.inr (f2_ne10 y.2 hv)
        · left; omega
    · have hA : a = shEmb m y := hsa.symm
      have hB : b = shEmb m x := hsb.symm
      subst hA
      subst hB
      have hxa : ((shEmb m x).1 : ℕ) = (x.1 : ℕ) + 1 := rfl
      have hxb : ((shEmb m y).1 : ℕ) = (y.1 : ℕ) + 1 := rfl
      have hx2 : ((shEmb m x).2 : ℕ) = (x.2 : ℕ) := rfl
      have hy2 : ((shEmb m y).2 : ℕ) = (y.2 : ℕ) := rfl
      have hx : (x.1 : ℕ) < m + 1 := x.1.isLt
      have hy : (y.1 : ℕ) < m + 1 := y.1.isLt
      refine ⟨⟨((lad_adj_emb m x y).mpr hadj).symm, ?_, ?_⟩,
        (disj_wE_iff m 1 (shEmb m y) (shEmb m x)).mpr ⟨?_, ?_, ?_, ?_⟩⟩
      · rcases f2_val_cases y.2 with hv | hv
        · refine Or.inr ⟨?_, ?_⟩
          · omega
          · omega
        · exact Or.inl (by omega)
      · rcases f2_val_cases x.2 with hv | hv
        · refine Or.inr ⟨?_, ?_⟩
          · omega
          · omega
        · exact Or.inl (by omega)
      · rcases f2_val_cases y.2 with hv | hv
        · exact Or.inr (f2_ne10 y.2 hv)
        · left; omega
      · rcases f2_val_cases y.2 with hv | hv
        · exact Or.inr (f2_ne10 y.2 hv)
        · left; omega
      · rcases f2_val_cases x.2 with hv | hv
        · exact Or.inr (f2_ne10 x.2 hv)
        · left; omega
      · rcases f2_val_cases x.2 with hv | hv
        · exact Or.inr (f2_ne10 x.2 hv)
        · left; omega

/-! ## §3 T4b 组装：p = a + 2c + a(n-2) -/

theorem f2_10 : (1 : Fin 2) ≠ (0 : Fin 2) := fun h => f2_01 h.symm

/-- wE1 与 wE0 跨行不相交（行 0 vs 行 1）。 -/
theorem wE_cross_disj (m : ℕ) : Disjoint (edgeSupp (wE m 1)) (edgeSupp (wE m 0)) := by
  have h := disjoint_esupp_iff ((⟨m+2, by omega⟩, (1 : Fin 2)) : Fin (m+3) × Fin 2)
    ((⟨0, by omega⟩, (1 : Fin 2)) : Fin (m+3) × Fin 2)
    ((⟨m+2, by omega⟩, (0 : Fin 2)) : Fin (m+3) × Fin 2)
    ((⟨0, by omega⟩, (0 : Fin 2)) : Fin (m+3) × Fin 2)
  refine h.mpr ⟨?_, ?_, ?_, ?_⟩
  · intro hcon; exact f2_10 (congrArg Prod.snd hcon)
  · intro hcon
    have hval := congrArg Fin.val (congrArg Prod.fst hcon)
    have h1 : ((⟨m+2, by omega⟩ : Fin (m+3)) : ℕ) = m + 2 := rfl
    have h2 : ((⟨0, by omega⟩ : Fin (m+3)) : ℕ) = 0 := rfl
    rw [h1, h2] at hval
    omega
  · intro hcon
    have hval := congrArg Fin.val (congrArg Prod.fst hcon)
    have h1 : ((⟨0, by omega⟩ : Fin (m+3)) : ℕ) = 0 := rfl
    have h2 : ((⟨m+2, by omega⟩ : Fin (m+3)) : ℕ) = m + 2 := rfl
    rw [h1, h2] at hval
    omega
  · intro hcon; exact f2_10 (congrArg Prod.snd hcon)

theorem wE1_not_mem_Cst (m : ℕ) : wE m 1 ∉ PrismR2.Cst (m+3) (m+3) := by
  intro h
  simp only [PrismR2.Cst, Finset.mem_filter] at h
  exact wE_not_mem_Eset m 1 h.1

/-- P \ {wE0} = Eset ∪ {wE1}。 -/
theorem pE_sdiff (m : ℕ) :
    (Eset (m+3) (m+3) ∪ {wE m 0, wE m 1}) \ {wE m 0} = Eset (m+3) (m+3) ∪ {wE m 1} := by
  ext f
  simp only [Finset.mem_sdiff, Finset.mem_union, Finset.mem_insert, Finset.mem_singleton]
  constructor
  · rintro ⟨(h | h | h), hne⟩
    · exact Or.inl h
    · exact absurd h hne
    · exact Or.inr h
  · rintro (h | h)
    · refine ⟨Or.inl h, ?_⟩
      intro hcon
      have hw : wE m 0 ∈ Eset (m+3) (m+3) := by rw [← hcon]; exact h
      exact wE_not_mem_Eset m 0 hw
    · refine ⟨Or.inr (Or.inr h), ?_⟩
      intro hcon
      exact wE_ne m (hcon.symm.trans h)

/-- (Eset ∪ {wE1}) \ {wE1} = Eset。 -/
theorem eset_sdiff (m : ℕ) :
    (Eset (m+3) (m+3) ∪ {wE m 1}) \ {wE m 1} = Eset (m+3) (m+3) := by
  ext f
  simp only [Finset.mem_sdiff, Finset.mem_union, Finset.mem_singleton]
  constructor
  · rintro ⟨(h | h), hne⟩
    · exact h
    · rw [h] at hne; exact absurd rfl hne
  · intro h
    refine ⟨Or.inl h, ?_⟩
    intro hcon
    have hw : wE m 1 ∈ Eset (m+3) (m+3) := by rw [← hcon]; exact h
    exact wE_not_mem_Eset m 1 hw

/-- (Cst ∪ {wE1}) \ {wE1} = Cst。 -/
theorem cst_sdiff (m : ℕ) :
    (PrismR2.Cst (m+3) (m+3) ∪ {wE m 1}) \ {wE m 1} = PrismR2.Cst (m+3) (m+3) := by
  ext f
  simp only [Finset.mem_sdiff, Finset.mem_union, Finset.mem_singleton]
  constructor
  · rintro ⟨(h | h), hne⟩
    · exact h
    · rw [h] at hne; exact absurd rfl hne
  · intro h
    refine ⟨Or.inl h, ?_⟩
    intro hcon
    have hw : wE m 1 ∈ PrismR2.Cst (m+3) (m+3) := by rw [← hcon]; exact h
    exact wE1_not_mem_Cst m hw

/-- (Eset ∪ {wE1}).filter (disj wE1) = Eset.filter (disj wE1)。 -/
theorem filter_u_wE1 (m : ℕ) :
    (Eset (m+3) (m+3) ∪ {wE m 1}).filter (fun f => Disjoint (edgeSupp f) (edgeSupp (wE m 1))) =
      (Eset (m+3) (m+3)).filter (fun f => Disjoint (edgeSupp f) (edgeSupp (wE m 1))) := by
  ext f
  simp only [Finset.mem_filter, Finset.mem_union, Finset.mem_singleton]
  constructor
  · rintro ⟨(h | h), hd⟩
    · exact ⟨h, hd⟩
    · rw [h] at hd; exact absurd hd (wE_nondisj m 1)
  · exact fun h => ⟨Or.inl h.1, h.2⟩

/-- (Cst ∪ {wE1}).filter (disj wE1) = Cst.filter (disj wE1)。 -/
theorem filter_c_wE1 (m : ℕ) :
    (PrismR2.Cst (m+3) (m+3) ∪ {wE m 1}).filter (fun f => Disjoint (edgeSupp f) (edgeSupp (wE m 1))) =
      (PrismR2.Cst (m+3) (m+3)).filter (fun f => Disjoint (edgeSupp f) (edgeSupp (wE m 1))) := by
  ext f
  simp only [Finset.mem_filter, Finset.mem_union, Finset.mem_singleton]
  constructor
  · rintro ⟨(h | h), hd⟩
    · exact ⟨h, hd⟩
    · rw [h] at hd; exact absurd hd (wE_nondisj m 1)
  · exact fun h => ⟨Or.inl h.1, h.2⟩

/-- {wE0, wE1}.filter (disj wE0) = {wE1}。 -/
theorem filter_pair_wE0 (m : ℕ) :
    ({wE m 0, wE m 1} : Finset (Sym2 (Fin (m+3) × Fin 2))).filter
        (fun f => Disjoint (edgeSupp f) (edgeSupp (wE m 0))) = {wE m 1} := by
  ext f
  simp only [Finset.mem_filter, Finset.mem_insert, Finset.mem_singleton]
  constructor
  · rintro ⟨(rfl | h), hd⟩
    · exact absurd hd (wE_nondisj m 0)
    · exact h
  · intro h
    rw [h]
    exact ⟨Or.inr rfl, wE_cross_disj m⟩

/-- P.filter (disj wE0) = Cst ∪ {wE1}（cell C2）。 -/
theorem pFilter (m : ℕ) :
    (Eset (m+3) (m+3) ∪ {wE m 0, wE m 1}).filter (fun f => Disjoint (edgeSupp f) (edgeSupp (wE m 0))) =
      PrismR2.Cst (m+3) (m+3) ∪ {wE m 1} := by
  rw [Finset.filter_union, cell_wE0, filter_pair_wE0]

/-- **T4b 第 2 步**：MC(Eset ∪ {wE1}) = a(n) + c(n)。 -/
theorem mc_union_wE1 (m : ℕ) :
    matchCount (Eset (m+3) (m+3) ∪ {wE m 1}) =
      hosoya (pathGraph (m+3) □ pathGraph 2) + matchCount (PrismR2.Cst (m+3) (m+3)) := by
  have hpart := matchCount_partition
    (E := Eset (m+3) (m+3) ∪ {wE m 1}) (e := wE m 1) (by simp) (wE_nondisj m 1)
  rw [eset_sdiff m, filter_u_wE1 m] at hpart
  have hE : matchCount (Eset (m+3) (m+3)) = hosoya (pathGraph (m+3) □ pathGraph 2) :=
    (transp (le_refl (m+3))).symm
  have hC : matchCount ((PrismR2.Cst (m+3) (m+3)).image
      (Sym2.map (swapv : Fin (m+3) × Fin 2 → Fin (m+3) × Fin 2))) =
    matchCount (PrismR2.Cst (m+3) (m+3)) :=
    (matchCount_map_inj (swapv : Fin (m+3) × Fin 2 → Fin (m+3) × Fin 2) swapv_inj _).symm
  rw [hpart, cell_wE1 m, hE, hC]

/-- **T4b 第 5 步**：MC(Cst ∪ {wE1}) = c(n) + a(n-2)。 -/
theorem mc_Cst_wE1 (m : ℕ) :
    matchCount (PrismR2.Cst (m+3) (m+3) ∪ {wE m 1}) =
      matchCount (PrismR2.Cst (m+3) (m+3)) + hosoya (pathGraph (m+1) □ pathGraph 2) := by
  have hpart := matchCount_partition
    (E := PrismR2.Cst (m+3) (m+3) ∪ {wE m 1}) (e := wE m 1) (by simp) (wE_nondisj m 1)
  rw [cst_sdiff m, filter_c_wE1 m] at hpart
  have hE : matchCount (Eset (m+1) (m+1)) = hosoya (pathGraph (m+1) □ pathGraph 2) :=
    (transp (le_refl (m+1))).symm
  have hC : matchCount ((Eset (m+1) (m+1)).image (Sym2.map (shEmb m))) =
      matchCount (Eset (m+1) (m+1)) :=
    (matchCount_map_inj (shEmb m) (shEmb_inj m) _).symm
  rw [hpart, cell_C3 m, hC, hE]

/-- **T4b 主分解**：p(n) = a(n) + 2c(n) + a(n-2)。 -/
theorem p_decomp (m : ℕ) :
    matchCount (Eset (m+3) (m+3) ∪ {wE m 0, wE m 1}) =
      hosoya (pathGraph (m+3) □ pathGraph 2) + 2 * matchCount (PrismR2.Cst (m+3) (m+3)) +
      hosoya (pathGraph (m+1) □ pathGraph 2) := by
  have hpart := matchCount_partition
    (E := Eset (m+3) (m+3) ∪ {wE m 0, wE m 1}) (e := wE m 0) (by simp) (wE_nondisj m 0)
  rw [pE_sdiff m, pFilter m, mc_union_wE1 m, mc_Cst_wE1 m] at hpart
  omega

/-! ## §4 ambient 不变性 + c-R4 + T4c 主 R4 -/

/-- Cst N m 的成员资格线性化（任意 ambient）。 -/
theorem cstNN' {N m : ℕ} (a b : Fin N × Fin 2) :
    s(a, b) ∈ PrismR2.Cst N m ↔ (pathGraph N □ pathGraph 2).Adj a b ∧
      (a.1 : ℕ) < m ∧ (b.1 : ℕ) < m ∧
      ((a.2 : ℕ) = 1 ∨ 1 ≤ (a.1 : ℕ) ∧ (a.1 : ℕ) + 2 ≤ m) ∧
      ((b.2 : ℕ) = 1 ∨ 1 ≤ (b.1 : ℕ) ∧ (b.1 : ℕ) + 2 ≤ m) := by
  simp only [PrismR2.Cst, Finset.mem_filter, PrismR2.filt2, eset_mem_mk, and_assoc]

/-- emb 保 Cst 成员（升 ambient）。 -/
theorem cst_emb_mem {N m : ℕ} (hk : m ≤ N) (x y : Fin m × Fin 2)
    (hadj : (pathGraph m □ pathGraph 2).Adj x y) (hxm : (x.1 : ℕ) < m) (hym : (y.1 : ℕ) < m)
    (hqx : (x.2 : ℕ) = 1 ∨ 1 ≤ (x.1 : ℕ) ∧ (x.1 : ℕ) + 2 ≤ m)
    (hqy : (y.2 : ℕ) = 1 ∨ 1 ≤ (y.1 : ℕ) ∧ (y.1 : ℕ) + 2 ≤ m) :
    s(emb hk x, emb hk y) ∈ PrismR2.Cst N m := by
  rw [cstNN']
  refine ⟨?_, hxm, hym, hqx, hqy⟩
  rw [lad_adj_iff] at hadj ⊢
  rcases hadj with ⟨hc, hr⟩ | ⟨hr, hc⟩
  · refine Or.inl ⟨?_, hr⟩
    apply Fin.val_injective
    have h1 : ((emb hk x).1 : ℕ) = (x.1 : ℕ) := rfl
    have h2 : ((emb hk y).1 : ℕ) = (y.1 : ℕ) := rfl
    rw [h1, h2]
    exact congrArg Fin.val hc
  · refine Or.inr ⟨hr, ?_⟩
    have h1 : ((emb hk x).1 : ℕ) = (x.1 : ℕ) := rfl
    have h2 : ((emb hk y).1 : ℕ) = (y.1 : ℕ) := rfl
    rcases hc with hc | hc <;> omega

/-- **ambient 不变（集合版）**：Cst N m = Cst m m 的 emb 像。 -/
theorem cst_ambient_image (N m : ℕ) (hk : m ≤ N) :
    PrismR2.Cst N m = (PrismR2.Cst m m).image (Sym2.map (emb hk)) := by
  ext e
  induction e using Sym2.ind with
  | _ a b =>
  rw [Finset.mem_image, cstNN']
  constructor
  · rintro ⟨hadj, hma, hmb, hqa, hqb⟩
    have hax : emb hk ((⟨(a.1 : ℕ), by omega⟩ : Fin m), a.2) = a := by
      refine Prod.ext_iff.mpr ⟨?_, rfl⟩
      exact Fin.val_injective rfl
    have hby : emb hk ((⟨(b.1 : ℕ), by omega⟩ : Fin m), b.2) = b := by
      refine Prod.ext_iff.mpr ⟨?_, rfl⟩
      exact Fin.val_injective rfl
    refine ⟨s(((⟨(a.1 : ℕ), by omega⟩ : Fin m), a.2), ((⟨(b.1 : ℕ), by omega⟩ : Fin m), b.2)),
      ?_, ?_⟩
    · rw [cstNN']
      refine ⟨?_, hma, hmb, hqa, hqb⟩
      rw [lad_adj_iff] at hadj ⊢
      rcases hadj with ⟨hc, hr⟩ | ⟨hr, hc⟩
      · refine Or.inl ⟨?_, hr⟩
        apply Fin.val_injective
        have h1 : ((⟨(a.1 : ℕ), by omega⟩ : Fin m) : ℕ) = (a.1 : ℕ) := rfl
        have h2 : ((⟨(b.1 : ℕ), by omega⟩ : Fin m) : ℕ) = (b.1 : ℕ) := rfl
        rw [h1, h2]
        exact congrArg Fin.val hc
      · refine Or.inr ⟨hr, ?_⟩
        have h1 : (((⟨(a.1 : ℕ), by omega⟩ : Fin m), a.2).1 : ℕ) = (a.1 : ℕ) := rfl
        have h2 : (((⟨(b.1 : ℕ), by omega⟩ : Fin m), b.2).1 : ℕ) = (b.1 : ℕ) := rfl
        rcases hc with hc | hc <;> omega
    · rw [Sym2.map_mk, hax, hby]
  · rintro ⟨e', he', hmap⟩
    induction e' using Sym2.ind with
    | _ x y =>
    rw [Sym2.map_mk] at hmap
    rw [cstNN'] at he'
    obtain ⟨hadj, hxm, hym, hqx, hqy⟩ := he'
    rw [Sym2.eq_iff] at hmap
    rcases hmap with ⟨h1, h2⟩ | ⟨h1, h2⟩
    · rw [← h1, ← h2]
      exact (cstNN' (emb hk x) (emb hk y)).mp (cst_emb_mem hk x y hadj hxm hym hqx hqy)
    · rw [← h2, ← h1]
      exact (cstNN' (emb hk y) (emb hk x)).mp
        (cst_emb_mem hk y x hadj.symm hym hxm hqy hqx)

/-- **ambient 不变（计数版）**：MC(Cst N m) = MC(Cst m m)。 -/
theorem cst_ambient (N m : ℕ) (hk : m ≤ N) :
    matchCount (PrismR2.Cst N m) = matchCount (PrismR2.Cst m m) := by
  rw [cst_ambient_image N m hk]
  exact (matchCount_map_inj (emb hk) (emb_inj hk) _).symm

/-- **c-R4**（ambient-free 口径，n ≥ 3）。 -/
theorem c_rec4' (n : ℕ) (hn : 3 ≤ n) :
    matchCount (PrismR2.Cst (n+4) (n+4)) + matchCount (PrismR2.Cst n n) =
      2 * matchCount (PrismR2.Cst (n+3) (n+3)) + 4 * matchCount (PrismR2.Cst (n+2) (n+2)) := by
  have h := PrismR2.c_rec4 (N := n+5) (n := n) hn (by omega)
  rw [cst_ambient (n+5) (n+4) (by omega), cst_ambient (n+5) n (by omega),
      cst_ambient (n+5) (n+3) (by omega), cst_ambient (n+5) (n+2) (by omega)] at h
  exact h

/-- **T4c**：棱柱主 R4。 -/
theorem hosoya_prism_rec4 (m : ℕ) :
    hosoya (prism (m+7)) + hosoya (prism (m+3)) =
      2 * hosoya (prism (m+6)) + 4 * hosoya (prism (m+5)) := by
  have hm4 : hosoya (prism (m+7)) = matchCount (Eset (m+7) (m+7) ∪ {wE (m+4) 0, wE (m+4) 1}) :=
    hosoya_prism_mc (m+4)
  have hm3 : hosoya (prism (m+6)) = matchCount (Eset (m+6) (m+6) ∪ {wE (m+3) 0, wE (m+3) 1}) :=
    hosoya_prism_mc (m+3)
  have hm2 : hosoya (prism (m+5)) = matchCount (Eset (m+5) (m+5) ∪ {wE (m+2) 0, wE (m+2) 1}) :=
    hosoya_prism_mc (m+2)
  have hm0 : hosoya (prism (m+3)) = matchCount (Eset (m+3) (m+3) ∪ {wE m 0, wE m 1}) :=
    hosoya_prism_mc m
  have pd4 : matchCount (Eset (m+7) (m+7) ∪ {wE (m+4) 0, wE (m+4) 1}) =
      hosoya (pathGraph (m+7) □ pathGraph 2) + 2 * matchCount (PrismR2.Cst (m+7) (m+7)) +
      hosoya (pathGraph (m+5) □ pathGraph 2) := p_decomp (m+4)
  have pd3 : matchCount (Eset (m+6) (m+6) ∪ {wE (m+3) 0, wE (m+3) 1}) =
      hosoya (pathGraph (m+6) □ pathGraph 2) + 2 * matchCount (PrismR2.Cst (m+6) (m+6)) +
      hosoya (pathGraph (m+4) □ pathGraph 2) := p_decomp (m+3)
  have pd2 : matchCount (Eset (m+5) (m+5) ∪ {wE (m+2) 0, wE (m+2) 1}) =
      hosoya (pathGraph (m+5) □ pathGraph 2) + 2 * matchCount (PrismR2.Cst (m+5) (m+5)) +
      hosoya (pathGraph (m+3) □ pathGraph 2) := p_decomp (m+2)
  have pd0 : matchCount (Eset (m+3) (m+3) ∪ {wE m 0, wE m 1}) =
      hosoya (pathGraph (m+3) □ pathGraph 2) + 2 * matchCount (PrismR2.Cst (m+3) (m+3)) +
      hosoya (pathGraph (m+1) □ pathGraph 2) := p_decomp m
  have ha0 : hosoya (pathGraph (m+4) □ pathGraph 2) + hosoya (pathGraph m □ pathGraph 2) =
      2 * hosoya (pathGraph (m+3) □ pathGraph 2) + 4 * hosoya (pathGraph (m+2) □ pathGraph 2) :=
    hosoya_ladder_rec4 m
  have ha1 : hosoya (pathGraph (m+5) □ pathGraph 2) + hosoya (pathGraph (m+1) □ pathGraph 2) =
      2 * hosoya (pathGraph (m+4) □ pathGraph 2) + 4 * hosoya (pathGraph (m+3) □ pathGraph 2) :=
    hosoya_ladder_rec4 (m+1)
  have ha2 : hosoya (pathGraph (m+6) □ pathGraph 2) + hosoya (pathGraph (m+2) □ pathGraph 2) =
      2 * hosoya (pathGraph (m+5) □ pathGraph 2) + 4 * hosoya (pathGraph (m+4) □ pathGraph 2) :=
    hosoya_ladder_rec4 (m+2)
  have ha3 : hosoya (pathGraph (m+7) □ pathGraph 2) + hosoya (pathGraph (m+3) □ pathGraph 2) =
      2 * hosoya (pathGraph (m+6) □ pathGraph 2) + 4 * hosoya (pathGraph (m+5) □ pathGraph 2) :=
    hosoya_ladder_rec4 (m+3)
  have hc : matchCount (PrismR2.Cst (m+7) (m+7)) + matchCount (PrismR2.Cst (m+3) (m+3)) =
      2 * matchCount (PrismR2.Cst (m+6) (m+6)) + 4 * matchCount (PrismR2.Cst (m+5) (m+5)) :=
    c_rec4' (m+3) (by omega)
  rw [hm4, hm0, hm3, hm2, pd4, pd0, pd3, pd2]
  omega

/-! ## §5 T4d rim 桥：整体旋转把 delRim 边集映成 Eset ∪ {wE1} -/

/-- 无环绕时 Fin 加一的 val。 -/
theorem fin_succ_val (m : ℕ) (i : Fin (m+3)) (h : (i : ℕ) + 1 < m + 3) :
    ((i + 1 : Fin (m+3)) : ℕ) = (i : ℕ) + 1 := by
  rw [Fin.val_add, Fin.val_one]
  exact Nat.mod_eq_of_lt h

/-- prismDelRim 删的那条 rim 边（第 0-1 列、行 0）。 -/
def e0 (m : ℕ) : Sym2 (Fin (m+3) × Fin 2) :=
  s((⟨0, by omega⟩, (0 : Fin 2)), (⟨1, by omega⟩, (0 : Fin 2)))

/-- 列旋转 −1（含环绕），行不变。 -/
def colw (m : ℕ) (i : Fin (m+3)) : Fin (m+3) :=
  if (i : ℕ) = 0 then ⟨m+2, by omega⟩ else ⟨(i : ℕ) - 1, by omega⟩

def rotw (m : ℕ) (v : Fin (m+3) × Fin 2) : Fin (m+3) × Fin 2 := (colw m v.1, v.2)

theorem rotw_row (m : ℕ) (v : Fin (m+3) × Fin 2) : (rotw m v).2 = v.2 := rfl

theorem colw_val_zero (m : ℕ) (i : Fin (m+3)) (h : (i : ℕ) = 0) : ((colw m i) : ℕ) = m + 2 := by
  simp only [colw, if_pos h]

theorem colw_val_pos (m : ℕ) (i : Fin (m+3)) (h : 1 ≤ (i : ℕ)) : ((colw m i) : ℕ) = (i : ℕ) - 1 := by
  have hne : ¬ ((i : ℕ) = 0) := by omega
  simp only [colw, if_neg hne]

theorem colw_zero (m : ℕ) : colw m (⟨0, by omega⟩ : Fin (m+3)) = ⟨m+2, by omega⟩ := rfl

theorem colw_one (m : ℕ) : colw m (⟨1, by omega⟩ : Fin (m+3)) = ⟨0, by omega⟩ := by
  have h1 : ((⟨1, by omega⟩ : Fin (m+3)) : ℕ) = 1 := rfl
  have hne : ¬ (((⟨1, by omega⟩ : Fin (m+3)) : ℕ) = 0) := by
    intro hcon
    rw [h1] at hcon
    omega
  simp only [colw, if_neg hne]

theorem colw_inj (m : ℕ) : Function.Injective (colw m) := by
  intro i j h
  by_cases hi : (i : ℕ) = 0
  · by_cases hj : (j : ℕ) = 0
    · exact Fin.val_injective (by show (i : ℕ) = (j : ℕ); omega)
    · exfalso
      have hv := congrArg Fin.val h
      rw [colw_val_zero m i hi, colw_val_pos m j (by omega)] at hv
      have his : (i : ℕ) < m + 3 := i.isLt
      omega
  · by_cases hj : (j : ℕ) = 0
    · exfalso
      have hv := congrArg Fin.val h
      rw [colw_val_pos m i (by omega), colw_val_zero m j hj] at hv
      have his : (i : ℕ) < m + 3 := i.isLt
      omega
    · exact Fin.val_injective (by
        have hv := congrArg Fin.val h
        rw [colw_val_pos m i (by omega), colw_val_pos m j (by omega)] at hv
        omega)

theorem rotw_inj (m : ℕ) : Function.Injective (rotw m) := by
  intro v v' h
  have hrow : (rotw m v).2 = (rotw m v').2 := congrArg Prod.snd h
  have hcol : (rotw m v).1 = (rotw m v').1 := congrArg Prod.fst h
  exact Prod.ext_iff.mpr ⟨colw_inj m hcol, hrow⟩

theorem rotv_rotw (m : ℕ) (v : Fin (m+3) × Fin 2) : rotv m (rotw m v) = v := by
  have hrow : (rotv m (rotw m v)).2 = v.2 := rfl
  have hc : (rotv m (rotw m v)).1 = v.1 := by
    show (colw m v.1) + 1 = v.1
    by_cases hv : (v.1 : ℕ) = 0
    · have hcw : ((colw m v.1) : ℕ) = m + 2 := colw_val_zero m v.1 hv
      apply Fin.val_injective
      rw [Fin.val_add, Fin.val_one, hcw]
      have hmod : (m + 2 + 1) % (m + 3) = 0 := by
        rw [show m + 2 + 1 = m + 3 from by omega]; exact Nat.mod_self _
      rw [hmod]
      omega
    · have hcw : ((colw m v.1) : ℕ) = (v.1 : ℕ) - 1 := colw_val_pos m v.1 (by omega)
      have his : (v.1 : ℕ) < m + 3 := v.1.isLt
      apply Fin.val_injective
      rw [Fin.val_add, Fin.val_one, hcw]
      rw [Nat.mod_eq_of_lt (by omega)]
      omega
  exact Prod.ext_iff.mpr ⟨hc, hrow⟩

theorem rotw_rotv (m : ℕ) (v : Fin (m+3) × Fin 2) : rotw m (rotv m v) = v := by
  have hrow : (rotw m (rotv m v)).2 = v.2 := rfl
  have hc : (rotw m (rotv m v)).1 = v.1 := by
    show colw m (v.1 + 1) = v.1
    by_cases hv : (v.1 : ℕ) = m + 2
    · have hvv : ((v.1 + 1 : Fin (m+3)) : ℕ) = 0 := by
        rw [Fin.val_add, Fin.val_one]
        have hmod : ((v.1 : ℕ) + 1) % (m + 3) = 0 := by
          rw [show (v.1 : ℕ) + 1 = m + 3 from by omega]; exact Nat.mod_self _
        rw [hmod]
      have hcw : ((colw m (v.1 + 1)) : ℕ) = m + 2 := colw_val_zero m (v.1 + 1) hvv
      apply Fin.val_injective
      rw [hcw]
      omega
    · have hvv : ((v.1 + 1 : Fin (m+3)) : ℕ) = (v.1 : ℕ) + 1 :=
        fin_succ_val m v.1 (by omega)
      have hpos : 1 ≤ ((v.1 + 1 : Fin (m+3)) : ℕ) := by omega
      have hcw : ((colw m (v.1 + 1)) : ℕ) = (v.1 : ℕ) := by
        rw [colw_val_pos m (v.1 + 1) hpos, hvv]
        omega
      apply Fin.val_injective
      rw [hcw]
  exact Prod.ext_iff.mpr ⟨hc, hrow⟩

/-- rotw 作用在 e0 上恰是 wE0。 -/
theorem rotw_e0 (m : ℕ) : Sym2.map (rotw m) (e0 m) = wE m 0 := by
  simp only [e0, Sym2.map_mk]
  have h1 : rotw m ((⟨0, by omega⟩, (0 : Fin 2)) : Fin (m+3) × Fin 2) =
      ((⟨m+2, by omega⟩, (0 : Fin 2)) : Fin (m+3) × Fin 2) := by
    refine Prod.ext_iff.mpr ⟨?_, rfl⟩
    show colw m (⟨0, by omega⟩ : Fin (m+3)) = _
    rw [colw_zero]
  have h2 : rotw m ((⟨1, by omega⟩, (0 : Fin 2)) : Fin (m+3) × Fin 2) =
      ((⟨0, by omega⟩, (0 : Fin 2)) : Fin (m+3) × Fin 2) := by
    refine Prod.ext_iff.mpr ⟨?_, rfl⟩
    show colw m (⟨1, by omega⟩ : Fin (m+3)) = _
    rw [colw_one]
  rw [h1, h2]
  simp only [wE]

/-- rotv 把 wE1 映到第 0-1 列行 1 的梯子边。 -/
theorem rotv_wE1 (m : ℕ) : Sym2.map (rotv m) (wE m 1) =
    s(((⟨0, by omega⟩, (1 : Fin 2)) : Fin (m+3) × Fin 2),
      ((⟨1, by omega⟩, (1 : Fin 2)) : Fin (m+3) × Fin 2)) := by
  simp only [wE, Sym2.map_mk]
  have h1 : rotv m ((⟨m+2, by omega⟩, (1 : Fin 2)) : Fin (m+3) × Fin 2) =
      ((⟨0, by omega⟩, (1 : Fin 2)) : Fin (m+3) × Fin 2) := by
    refine Prod.ext_iff.mpr ⟨?_, rfl⟩
    apply Fin.val_injective
    show (((⟨m+2, by omega⟩ : Fin (m+3)) + 1 : Fin (m+3)) : ℕ) = 0
    rw [Fin.val_add, Fin.val_one, top_val m]
    have hmod : (m + 2 + 1) % (m + 3) = 0 := by
      rw [show m + 2 + 1 = m + 3 from by omega]; exact Nat.mod_self _
    rw [hmod]
  have h2 : rotv m ((⟨0, by omega⟩, (1 : Fin 2)) : Fin (m+3) × Fin 2) =
      ((⟨1, by omega⟩, (1 : Fin 2)) : Fin (m+3) × Fin 2) := by
    refine Prod.ext_iff.mpr ⟨?_, rfl⟩
    apply Fin.val_injective
    show (((⟨0, by omega⟩ : Fin (m+3)) + 1 : Fin (m+3)) : ℕ) = 1
    rw [Fin.val_add, Fin.val_one, zero_val m]
    have hmod : (0 + 1) % (m + 3) = 1 := Nat.mod_eq_of_lt (by omega)
    rw [hmod]
  rw [h1, h2]

/-- 梯子邻接在 rotv 下的像仍是棱柱邻接。 -/
theorem prism_adj_rotv (m : ℕ) (x y : Fin (m+3) × Fin 2)
    (hadj : (pathGraph (m+3) □ pathGraph 2).Adj x y) :
    (prism (m+3)).Adj (rotv m x) (rotv m y) := by
  show (cycleGraph (m+3) □ pathGraph 2).Adj (rotv m x) (rotv m y)
  rw [boxProd_adj]
  have his : (x.1 : ℕ) < m + 3 := x.1.isLt
  rw [lad_adj_iff] at hadj
  rcases hadj with ⟨hc, hr⟩ | ⟨hr, hc⟩
  · refine Or.inr ⟨(p2adj _ _).mpr hr, ?_⟩
    show (x.1 + 1 : Fin (m+3)) = (y.1 + 1 : Fin (m+3))
    exact congrArg (fun i : Fin (m+3) => i + 1) hc
  · refine Or.inl ⟨?_, hr⟩
    rw [cycleGraph_adj_iff_succ]
    rcases hc with hc | hc
    · refine Or.inl ?_
      have hfin : x.1 + 1 = y.1 := by
        apply Fin.val_injective
        rw [Fin.val_add, Fin.val_one]
        rw [show ((x.1 : ℕ) + 1) % (m + 3) = (x.1 : ℕ) + 1 from Nat.mod_eq_of_lt (by omega)]
        exact hc
      show (y.1 + 1 : Fin (m+3)) = ((x.1 + 1 : Fin (m+3)) + 1)
      exact (congrArg (fun i : Fin (m+3) => i + 1) hfin).symm
    · refine Or.inr ?_
      have hfin : y.1 + 1 = x.1 := by
        apply Fin.val_injective
        rw [Fin.val_add, Fin.val_one]
        rw [show ((y.1 : ℕ) + 1) % (m + 3) = (y.1 : ℕ) + 1 from Nat.mod_eq_of_lt (by omega)]
        exact hc
      show (x.1 + 1 : Fin (m+3)) = ((y.1 + 1 : Fin (m+3)) + 1)
      exact (congrArg (fun i : Fin (m+3) => i + 1) hfin).symm

/-- 第 0-1 列行 1 的梯子边是棱柱边。 -/
theorem adj01 (m : ℕ) :
    (prism (m+3)).Adj ((⟨0, by omega⟩, (1 : Fin 2)) : Fin (m+3) × Fin 2)
      ((⟨1, by omega⟩, (1 : Fin 2)) : Fin (m+3) × Fin 2) := by
  show (cycleGraph (m+3) □ pathGraph 2).Adj _ _
  refine boxProd_adj.mpr (Or.inl ⟨?_, rfl⟩)
  rw [cycleGraph_adj_iff_succ]
  refine Or.inl ?_
  apply Fin.val_injective
  have hA : ((⟨1, by omega⟩ : Fin (m+3)) : ℕ) = 1 := rfl
  have hB : ((⟨0, by omega⟩ : Fin (m+3)) : ℕ) = 0 := rfl
  rw [Fin.val_add, Fin.val_one, hA, hB]
  exact (Nat.mod_eq_of_lt (by omega)).symm

/-- 行 1 的边不是 e0。 -/
theorem ne01 (m : ℕ) :
    s(((⟨0, by omega⟩, (1 : Fin 2)) : Fin (m+3) × Fin 2),
      ((⟨1, by omega⟩, (1 : Fin 2)) : Fin (m+3) × Fin 2)) ≠ e0 m := by
  intro hcon
  simp only [e0] at hcon
  rw [Sym2.eq_iff] at hcon
  rcases hcon with ⟨h1, h2⟩ | ⟨h1, h2⟩
  · exact f2_10 (congrArg Prod.snd h1)
  · exact f2_10 (congrArg Prod.snd h2)

theorem sym2_swap (V : Type*) (x y : V) : s(x, y) = s(y, x) :=
  (Sym2.eq_iff.mpr (Or.inr ⟨rfl, rfl⟩))

/-- **cell C4（旋转版）**：(Eset ∪ {wE1}) 的 rotv 像 = prismDelRim 边集。 -/
theorem cell_C4 (m : ℕ) :
    (Eset (m+3) (m+3) ∪ {wE m 1}).image (Sym2.map (rotv m)) =
      (prism (m+3)).edgeFinset \ {e0 m} := by
  ext e
  induction e using Sym2.ind with
  | _ a b =>
  simp only [Finset.mem_image, Finset.mem_sdiff, Finset.mem_union, Finset.mem_singleton,
    SimpleGraph.mem_edgeFinset, SimpleGraph.mem_edgeSet]
  constructor
  · rintro ⟨e', hmem, hmap⟩
    induction e' using Sym2.ind with
    | _ x y =>
    rw [Sym2.map_mk] at hmap
    rcases hmem with hE | hw
    · have hadj : (prism (m+3)).Adj (rotv m x) (rotv m y) :=
        prism_adj_rotv m x y ((esetNN m x y).mp hE)
      have hne : s(rotv m x, rotv m y) ≠ e0 m := by
        intro hcon
        have hc2 : Sym2.map (rotw m) (s(rotv m x, rotv m y)) = Sym2.map (rotw m) (e0 m) :=
          congrArg (Sym2.map (rotw m)) hcon
        rw [Sym2.map_mk, rotw_rotv, rotw_rotv, rotw_e0] at hc2
        exact wE_not_mem_Eset m 0 (by rw [← hc2]; exact hE)
      rcases (Sym2.eq_iff.mp hmap) with ⟨h1, h2⟩ | ⟨h1, h2⟩
      · rw [← h1, ← h2]
        exact ⟨hadj, hne⟩
      · rw [← h2, ← h1]
        exact ⟨hadj.symm,
          fun hcon => hne ((sym2_swap _ (rotv m x) (rotv m y)).trans hcon)⟩
    · have hmma : Sym2.map (rotv m) (s(x, y)) = s(rotv m x, rotv m y) := Sym2.map_mk (rotv m) x y
      rw [hw, rotv_wE1] at hmma
      have hab := hmma.trans hmap
      rw [Sym2.eq_iff] at hab
      rcases hab with ⟨h1, h2⟩ | ⟨h1, h2⟩
      · rw [← h1, ← h2]
        exact ⟨adj01 m, ne01 m⟩
      · rw [← h1, ← h2]
        exact ⟨(adj01 m).symm, fun hcon =>
          ne01 m ((sym2_swap _
            ((⟨0, by omega⟩, (1 : Fin 2)) : Fin (m+3) × Fin 2)
            ((⟨1, by omega⟩, (1 : Fin 2)) : Fin (m+3) × Fin 2)).trans hcon)⟩
  · rintro ⟨hadj, hne⟩
    refine ⟨s(rotw m a, rotw m b), ?_, ?_⟩
    · rw [show (prism (m+3)).Adj a b = ((cycleGraph (m+3) □ pathGraph 2).Adj a b) from rfl] at hadj
      rw [boxProd_adj] at hadj
      rcases hadj with ⟨hcyc, hrow⟩ | ⟨hp2, hcol⟩
      · rw [cycleGraph_adj_iff_succ] at hcyc
        rcases f2_val_cases a.2 with hv | hv
        · -- 行 0：a.2 = 0
          by_cases ha0 : (a.1 : ℕ) = 0
          · rcases hcyc with hb | hb
            · -- b.1 = a.1+1 = ⟨1⟩ → 与 ≠e0 矛盾
              exfalso
              apply hne
              have hvt : a.1 = (⟨0, by omega⟩ : Fin (m+3)) :=
                Fin.val_injective (by show (a.1 : ℕ) = 0; omega)
              have hbv : (b.1 : ℕ) = 1 := by
                rcases (fin_add_one_val m a.1 b.1).mp hb with h | ⟨h, h2⟩
                · omega
                · omega
              have hbt : b.1 = (⟨1, by omega⟩ : Fin (m+3)) :=
                Fin.val_injective (by show (b.1 : ℕ) = 1; omega)
              have hrowb : b.2 = (0 : Fin 2) := by
                rw [← hrow]; exact f2_val0 a.2 hv
              have hat : a = ((⟨0, by omega⟩, (0 : Fin 2)) : Fin (m+3) × Fin 2) :=
                Prod.ext_iff.mpr ⟨hvt, f2_val0 a.2 hv⟩
              have hbt2 : b = ((⟨1, by omega⟩, (0 : Fin 2)) : Fin (m+3) × Fin 2) :=
                Prod.ext_iff.mpr ⟨hbt, hrowb⟩
              rw [hat, hbt2]
              rfl
            · -- a.1 = b.1+1，b.1.val = m+2：梯子边
              refine Or.inl ?_
              rw [esetNN, lad_adj_iff]
              refine Or.inr ⟨hrow, ?_⟩
              have hbv : (b.1 : ℕ) = m + 2 := by
                rcases (fin_add_one_val m b.1 a.1).mp hb with h | ⟨h, h2⟩
                · omega
                · omega
              have hca : ((rotw m a).1 : ℕ) = m + 2 := colw_val_zero m a.1 ha0
              have hcb : ((rotw m b).1 : ℕ) = (b.1 : ℕ) - 1 := colw_val_pos m b.1 (by omega)
              rcases f2_val_cases b.2 with hvb | hvb
              · right; omega
              · right; omega
          · by_cases hb0 : (b.1 : ℕ) = 0
            · rcases hcyc with hb | hb
              · -- b.1 = a.1+1（a.1.val = m+2 环绕）：梯子边
                refine Or.inl ?_
                rw [esetNN, lad_adj_iff]
                refine Or.inr ⟨hrow, ?_⟩
                have hav : (a.1 : ℕ) = m + 2 := by
                  rcases (fin_add_one_val m a.1 b.1).mp hb with h | ⟨h, h2⟩
                  · omega
                  · omega
                have hca : ((rotw m a).1 : ℕ) = (a.1 : ℕ) - 1 := colw_val_pos m a.1 (by omega)
                have hcb : ((rotw m b).1 : ℕ) = m + 2 := colw_val_zero m b.1 hb0
                rcases f2_val_cases b.2 with hvb | hvb
                · left; omega
                · left; omega
              · -- a.1 = b.1+1 = ⟨1⟩ → 与 ≠e0 矛盾
                exfalso
                apply hne
                have hbt : b.1 = (⟨0, by omega⟩ : Fin (m+3)) :=
                  Fin.val_injective (by show (b.1 : ℕ) = 0; omega)
                have hav : (a.1 : ℕ) = 1 := by
                  rcases (fin_add_one_val m b.1 a.1).mp hb with h | ⟨h, h2⟩
                  · omega
                  · omega
                have hat : a.1 = (⟨1, by omega⟩ : Fin (m+3)) :=
                  Fin.val_injective (by show (a.1 : ℕ) = 1; omega)
                have hrowa : a.2 = (0 : Fin 2) := f2_val0 a.2 hv
                have hrowb : b.2 = (0 : Fin 2) := by
                  rw [← hrow]; exact hrowa
                have hat2 : a = ((⟨1, by omega⟩, (0 : Fin 2)) : Fin (m+3) × Fin 2) :=
                  Prod.ext_iff.mpr ⟨hat, hrowa⟩
                have hbt2 : b = ((⟨0, by omega⟩, (0 : Fin 2)) : Fin (m+3) × Fin 2) :=
                  Prod.ext_iff.mpr ⟨hbt, hrowb⟩
                rw [hat2, hbt2]
                exact sym2_swap _ ((⟨1, by omega⟩, (0 : Fin 2)) : Fin (m+3) × Fin 2)
                  ((⟨0, by omega⟩, (0 : Fin 2)) : Fin (m+3) × Fin 2)
            · -- 两侧 val ≥ 1：colw = val−1，梯子邻接
              refine Or.inl ?_
              rw [esetNN, lad_adj_iff]
              refine Or.inr ⟨hrow, ?_⟩
              have hca : ((rotw m a).1 : ℕ) = (a.1 : ℕ) - 1 := colw_val_pos m a.1 (by omega)
              have hcb : ((rotw m b).1 : ℕ) = (b.1 : ℕ) - 1 := colw_val_pos m b.1 (by omega)
              rcases hcyc with hb | hb
              · rcases (fin_add_one_val m a.1 b.1).mp hb with h | ⟨h, h2⟩
                · left; omega
                · omega
              · rcases (fin_add_one_val m b.1 a.1).mp hb with h | ⟨h, h2⟩
                · right; omega
                · omega
        · -- 行 1：a.2 = 1
          by_cases ha0 : (a.1 : ℕ) = 0
          · rcases hcyc with hb | hb
            · -- b.1 = ⟨1⟩：像恰是 wE1
              refine Or.inr ?_
              have hvt : a.1 = (⟨0, by omega⟩ : Fin (m+3)) :=
                Fin.val_injective (by show (a.1 : ℕ) = 0; omega)
              have hbv : (b.1 : ℕ) = 1 := by
                rcases (fin_add_one_val m a.1 b.1).mp hb with h | ⟨h, h2⟩
                · omega
                · omega
              have hbt : b.1 = (⟨1, by omega⟩ : Fin (m+3)) :=
                Fin.val_injective (by show (b.1 : ℕ) = 1; omega)
              have hrowa : a.2 = (1 : Fin 2) := f2_val1 a.2 hv
              have hrowb : b.2 = (1 : Fin 2) := by rw [← hrow]; exact hrowa
              have ha : rotw m a = ((⟨m+2, by omega⟩, (1 : Fin 2)) : Fin (m+3) × Fin 2) := by
                refine Prod.ext_iff.mpr ⟨?_, hrowa⟩
                show colw m a.1 = _
                rw [hvt, colw_zero]
              have hb : rotw m b = ((⟨0, by omega⟩, (1 : Fin 2)) : Fin (m+3) × Fin 2) := by
                refine Prod.ext_iff.mpr ⟨?_, hrowb⟩
                show colw m b.1 = _
                rw [hbt, colw_one]
              rw [ha, hb]
              simp only [wE]
            · -- a.1 = b.1+1，b.1.val = m+2：梯子边
              refine Or.inl ?_
              rw [esetNN, lad_adj_iff]
              refine Or.inr ⟨hrow, ?_⟩
              have hbv : (b.1 : ℕ) = m + 2 := by
                rcases (fin_add_one_val m b.1 a.1).mp hb with h | ⟨h, h2⟩
                · omega
                · omega
              have hca : ((rotw m a).1 : ℕ) = m + 2 := colw_val_zero m a.1 ha0
              have hcb : ((rotw m b).1 : ℕ) = (b.1 : ℕ) - 1 := colw_val_pos m b.1 (by omega)
              rcases f2_val_cases b.2 with hvb | hvb
              · right; omega
              · right; omega
          · by_cases hb0 : (b.1 : ℕ) = 0
            · rcases hcyc with hb | hb
              · -- b.1 = a.1+1（a.1.val = m+2 环绕）：梯子边
                refine Or.inl ?_
                rw [esetNN, lad_adj_iff]
                refine Or.inr ⟨hrow, ?_⟩
                have hav : (a.1 : ℕ) = m + 2 := by
                  rcases (fin_add_one_val m a.1 b.1).mp hb with h | ⟨h, h2⟩
                  · omega
                  · omega
                have hca : ((rotw m a).1 : ℕ) = (a.1 : ℕ) - 1 := colw_val_pos m a.1 (by omega)
                have hcb : ((rotw m b).1 : ℕ) = m + 2 := colw_val_zero m b.1 hb0
                rcases f2_val_cases b.2 with hvb | hvb
                · left; omega
                · left; omega
              · -- a.1 = b.1+1 = ⟨1⟩：像恰是 wE1（对称）
                refine Or.inr ?_
                have hbt : b.1 = (⟨0, by omega⟩ : Fin (m+3)) :=
                  Fin.val_injective (by show (b.1 : ℕ) = 0; omega)
                have hav : (a.1 : ℕ) = 1 := by
                  rcases (fin_add_one_val m b.1 a.1).mp hb with h | ⟨h, h2⟩
                  · omega
                  · omega
                have hat : a.1 = (⟨1, by omega⟩ : Fin (m+3)) :=
                  Fin.val_injective (by show (a.1 : ℕ) = 1; omega)
                have hrowa : a.2 = (1 : Fin 2) := f2_val1 a.2 hv
                have hrowb : b.2 = (1 : Fin 2) := by rw [← hrow]; exact hrowa
                have ha : rotw m a = ((⟨0, by omega⟩, (1 : Fin 2)) : Fin (m+3) × Fin 2) := by
                  refine Prod.ext_iff.mpr ⟨?_, hrowa⟩
                  show colw m a.1 = _
                  rw [hat, colw_one]
                have hb : rotw m b = ((⟨m+2, by omega⟩, (1 : Fin 2)) : Fin (m+3) × Fin 2) := by
                  refine Prod.ext_iff.mpr ⟨?_, hrowb⟩
                  show colw m b.1 = _
                  rw [hbt, colw_zero]
                rw [ha, hb]
                simp only [wE]
                exact sym2_swap _ _ _
            · -- 两侧 val ≥ 1：梯子邻接
              refine Or.inl ?_
              rw [esetNN, lad_adj_iff]
              refine Or.inr ⟨hrow, ?_⟩
              have hca : ((rotw m a).1 : ℕ) = (a.1 : ℕ) - 1 := colw_val_pos m a.1 (by omega)
              have hcb : ((rotw m b).1 : ℕ) = (b.1 : ℕ) - 1 := colw_val_pos m b.1 (by omega)
              rcases hcyc with hb | hb
              · rcases (fin_add_one_val m a.1 b.1).mp hb with h | ⟨h, h2⟩
                · left; omega
                · omega
              · rcases (fin_add_one_val m b.1 a.1).mp hb with h | ⟨h, h2⟩
                · right; omega
                · omega
      · -- 竖档：行翻转、列不变
        refine Or.inl ?_
        rw [esetNN, lad_adj_iff]
        refine Or.inl ⟨congrArg (colw m) hcol, (p2adj _ _).mp hp2⟩
    · rw [Sym2.map_mk, rotv_rotw, rotv_rotw]

/-- prismDelRim 的边集 = 棱柱边集 \ {e0}。 -/
theorem delrim_edges (m : ℕ) (hm : 3 ≤ m + 3) :
    (prismDelRim (m+3) hm).edgeFinset = (prism (m+3)).edgeFinset \ {e0 m} := by
  ext f
  induction f using Sym2.ind with
  | _ x y =>
    simp only [Finset.mem_sdiff, Finset.mem_singleton]
    rw [SimpleGraph.mem_edgeFinset, SimpleGraph.mem_edgeSet,
      SimpleGraph.mem_edgeFinset, SimpleGraph.mem_edgeSet]
    simp only [prismDelRim]
    rw [SimpleGraph.deleteEdges_adj, Set.mem_singleton_iff]
    simp only [e0]
    rfl

/-- **T4d**：d_rim(n) = a(n) + c(n)。 -/
theorem delrim_eq (m : ℕ) (hm : 3 ≤ m + 3) :
    hosoya (prismDelRim (m+3) hm) =
      hosoya (pathGraph (m+3) □ pathGraph 2) + matchCount (PrismR2.Cst (m+3) (m+3)) := by
  rw [hosoya_matchCount, delrim_edges m hm, ← cell_C4]
  rw [← (matchCount_map_inj (rotv m) (rotv_inj m) (Eset (m+3) (m+3) ∪ {wE m 1}))]
  exact mc_union_wE1 m

/-! ## §6 T4e spoke 桥 + T4f 初值/强归纳 + 主桥收口

探针 provenance：T4e = 子代理 agent-15/18（双批同参，探针 probe-t4e.lean EXIT=0）；
初值 = agent-16/19（probe-t4f-init.lean，c6 经库件 T1–T7 递推链，decide 超界留档）；
强归纳/R4 = agent-17/20（probe-t4f-ind.lean）。组装 = 主代理 scaffold。
-/

set_option maxRecDepth 40000
set_option maxHeartbeats 1000000

-- ===== T4e spoke 桥 =====

/-- **T4e**：d_spoke 分割恒等式——删 spoke 的棱柱匹配数 = 完整棱柱 − (m+2) 梯子。 -/
theorem delspoke_eq (m : ℕ) (hm : 3 ≤ m + 3) :
    hosoya (prism (m+3)) =
      hosoya (prismDelSpoke (m+3) hm) + hosoya (pathGraph (m+2) □ pathGraph 2) := by
  have h := hosoya_prism_split_spoke (m+3) hm
  have h2 : hosoya (SimpleGraph.induce
      (Set.univ \ {((⟨0, by omega⟩, (0 : Fin 2))), ((⟨0, by omega⟩, (1 : Fin 2)))} :
        Set (Fin (m+3) × Fin 2)) (prism (m+3))) =
      hosoya (pathGraph (m+2) □ pathGraph 2) :=
    hosoya_iso _ _ (prism_delSpokeEnds_induce_iso_ladder (m+3) hm)
  rw [h2] at h
  exact h

-- ===== T4f 初值：a(4..6) = 71,228,733 与 c(3..6) = 4,15,46,150 =====

theorem a4 : hosoya (pathGraph 4 □ pathGraph 2) = 71 := by
  -- rec3 (m=1)：a(4) + a(1) = 3·a(3) + a(2)
  have hr := hosoya_ladder_rec3 1
  have h4 : hosoya (pathGraph (1 + 3) □ pathGraph 2) = hosoya (pathGraph 4 □ pathGraph 2) := rfl
  have h3 : hosoya (pathGraph (1 + 2) □ pathGraph 2) = hosoya (pathGraph 3 □ pathGraph 2) := rfl
  have h2 : hosoya (pathGraph (1 + 1) □ pathGraph 2) = hosoya (pathGraph 2 □ pathGraph 2) := rfl
  rw [h4, h3, h2] at hr
  have e1 := hosoya_lad1
  have e2 := hosoya_lad2
  have e3 := hosoya_lad3
  omega

theorem a5 : hosoya (pathGraph 5 □ pathGraph 2) = 228 := by
  -- rec3 (m=2)：a(5) + a(2) = 3·a(4) + a(3)
  have hr := hosoya_ladder_rec3 2
  have h5 : hosoya (pathGraph (2 + 3) □ pathGraph 2) = hosoya (pathGraph 5 □ pathGraph 2) := rfl
  have h4 : hosoya (pathGraph (2 + 2) □ pathGraph 2) = hosoya (pathGraph 4 □ pathGraph 2) := rfl
  have h3 : hosoya (pathGraph (2 + 1) □ pathGraph 2) = hosoya (pathGraph 3 □ pathGraph 2) := rfl
  rw [h5, h4, h3] at hr
  have e2 := hosoya_lad2
  have e3 := hosoya_lad3
  have e4 := a4
  omega

theorem a6 : hosoya (pathGraph 6 □ pathGraph 2) = 733 := by
  -- rec3 (m=3)：a(6) + a(3) = 3·a(5) + a(4)
  have hr := hosoya_ladder_rec3 3
  have h6 : hosoya (pathGraph (3 + 3) □ pathGraph 2) = hosoya (pathGraph 6 □ pathGraph 2) := rfl
  have h5 : hosoya (pathGraph (3 + 2) □ pathGraph 2) = hosoya (pathGraph 5 □ pathGraph 2) := rfl
  have h4 : hosoya (pathGraph (3 + 1) □ pathGraph 2) = hosoya (pathGraph 4 □ pathGraph 2) := rfl
  rw [h6, h5, h4] at hr
  have e3 := hosoya_lad3
  have e4 := a4
  have e5 := a5
  omega

theorem c3 : matchCount (PrismR2.Cst 3 3) = 4 := by decide

theorem c4 : matchCount (PrismR2.Cst 4 4) = 15 := by decide

theorem c5 : matchCount (PrismR2.Cst 5 5) = 46 := by decide

/-- c6 递推链所需的 4 个小实例（m ≤ 2 或与 c4 同规模的 decide）。 -/
private theorem c6_small :
    matchCount (PrismR2.V3 6 2) = 14 ∧
    matchCount (PrismR2.V5 6 2) = 10 ∧
    matchCount (PrismR2.V7 6 2) = 8 ∧
    matchCount (PrismR2.Cst 6 (3 + 1)) = 15 := by
  refine ⟨?_, ?_, ?_, ?_⟩ <;>
    (try decide) <;>
    (try decide) <;>
    (try decide) <;>
    (try decide)

/-- c(6) = 150：直接 decide 超界（2^21 幂集枚举，whnf 心跳超时留档），
改走 PrismHosoya 库件 T1–T7 递推实例链 + omega 线性消元。 -/
theorem c6 : matchCount (PrismR2.Cst 6 6) = 150 := by
  obtain ⟨e32, e42, e72, c6c4⟩ := c6_small
  set i3 : Fin 6 := ⟨3, by omega⟩ with hi3
  set i4 : Fin 6 := ⟨4, by omega⟩ with hi4
  -- V5r(3) = V3(2) + V7(2) = 14 + 8 = 22（T1 m=2）
  have s3 : matchCount (PrismR2.V5r 6 3 i3) = matchCount (PrismR2.V5r 6 (2 + 1) i3) := rfl
  have A1 := PrismR2.T1 (N := 6) (m := 2) i3 (by simp [hi3]) (by omega) (by omega)
  rw [← s3] at A1
  have hV5r3 : matchCount (PrismR2.V5r 6 3 i3) = 22 := by omega
  -- V5(3) = V5r(3) + V5(2) = 22 + 10 = 32（T7 m=2）
  have s5a : matchCount (PrismR2.V5 6 3) = matchCount (PrismR2.V5 6 (2 + 1)) := rfl
  have s5b : matchCount (PrismR2.V5r 6 3 i3) = matchCount (PrismR2.V5r 6 (2 + 1) i3) := rfl
  have A2 := PrismR2.T7 (N := 6) (m := 2) i3 (by simp [hi3]) (by omega) (by omega)
  rw [← s5a, ← s5b] at A2
  have hV53 : matchCount (PrismR2.V5 6 3) = 32 := by omega
  -- V3r(3) = V5r(3) + Cst(4) = 22 + 15 = 37（T4 m=3）
  have A3 := PrismR2.T4 (N := 6) (m := 3) i3 (by simp [hi3]) (by omega) (by omega)
  have hV3r3 : matchCount (PrismR2.V3r 6 3 i3) = 37 := by omega
  -- V3(3) = V3r(3) + V5(2) = 37 + 10 = 47（T3 m=3）
  have s7 : matchCount (PrismR2.V5 6 2) = matchCount (PrismR2.V5 6 (3 - 1)) := rfl
  have A4 := PrismR2.T3 (N := 6) (m := 3) i3 (by simp [hi3]) (by omega) (by omega)
  rw [← s7] at A4
  have hV33 : matchCount (PrismR2.V3 6 3) = 47 := by omega
  -- V7(3) = V3(2) + V5(2) = 14 + 10 = 24（T6 m=3）
  have s8a : matchCount (PrismR2.V3 6 2) = matchCount (PrismR2.V3 6 (3 - 1)) := rfl
  have A5 := PrismR2.T6 (N := 6) (m := 3) (by omega) (by omega)
  rw [← s8a, ← s7] at A5
  have hV73 : matchCount (PrismR2.V7 6 3) = 24 := by omega
  -- V5r(4) = V3(3) + V7(3) = 47 + 24 = 71（T1 m=3）
  have s9 : matchCount (PrismR2.V5r 6 4 i4) = matchCount (PrismR2.V5r 6 (3 + 1) i4) := rfl
  have A6 := PrismR2.T1 (N := 6) (m := 3) i4 (by simp [hi4]) (by omega) (by omega)
  rw [← s9] at A6
  have hV5r4 : matchCount (PrismR2.V5r 6 4 i4) = 71 := by omega
  -- V5(4) = V5r(4) + V5(3) = 71 + 32 = 103（T5 m=4）
  have s10 : matchCount (PrismR2.V5 6 3) = matchCount (PrismR2.V5 6 (4 - 1)) := rfl
  have A7 := PrismR2.T5 (N := 6) (m := 4) i4 (by simp [hi4]) (by omega) (by omega)
  rw [← s10] at A7
  have hV54 : matchCount (PrismR2.V5 6 4) = 103 := by omega
  -- Cst(6) = V5(4) + V3(3) = 103 + 47 = 150（T2 m=4）
  have s11a : matchCount (PrismR2.Cst 6 6) = matchCount (PrismR2.Cst 6 (4 + 2)) := rfl
  have s11b : matchCount (PrismR2.V3 6 3) = matchCount (PrismR2.V3 6 (4 - 1)) := rfl
  have A8 := PrismR2.T2 (N := 6) (m := 4) (by omega) (by omega)
  rw [← s11a, ← s11b] at A8
  omega

-- ===== T4f 强归纳 + R4 代数（纯 ℤ，探针逐字） =====

theorem aspoke_strong (D : ℕ → ℤ) (h0 : D 0 = 25) (h1 : D 1 = 86) (h2 : D 2 = 271) (h3 : D 3 = 876)
    (hr : ∀ k : ℕ, D (k+4) + D k = 2 * D (k+3) + 4 * D (k+2)) : ∀ m : ℕ, D m = aspoke m := by
  intro m
  induction m using Nat.strongRecOn with
  | ind n ih =>
    rcases Nat.lt_or_ge n 4 with hlt | hge
    · obtain rfl | rfl | rfl | rfl :=
        (show n = 0 ∨ n = 1 ∨ n = 2 ∨ n = 3 by omega)
      · exact h0
      · exact h1
      · exact h2
      · exact h3
    · obtain ⟨k, rfl⟩ : ∃ k, n = k + 4 := ⟨n - 4, by omega⟩
      have hrk := hr k
      have e3 := ih (k + 3) (by omega)
      have e2 := ih (k + 2) (by omega)
      have e0 := ih k (by omega)
      simp only [aspoke]
      omega

theorem arim_strong (D : ℕ → ℤ) (h0 : D 0 = 26) (h1 : D 1 = 86) (h2 : D 2 = 274) (h3 : D 3 = 883)
    (hr : ∀ k : ℕ, D (k+4) + D k = 2 * D (k+3) + 4 * D (k+2)) : ∀ m : ℕ, D m = arim m := by
  intro m
  induction m using Nat.strongRecOn with
  | ind n ih =>
    rcases Nat.lt_or_ge n 4 with hlt | hge
    · obtain rfl | rfl | rfl | rfl :=
        (show n = 0 ∨ n = 1 ∨ n = 2 ∨ n = 3 by omega)
      · exact h0
      · exact h1
      · exact h2
      · exact h3
    · obtain ⟨k, rfl⟩ : ∃ k, n = k + 4 := ⟨n - 4, by omega⟩
      have hrk := hr k
      have e3 := ih (k + 3) (by omega)
      have e2 := ih (k + 2) (by omega)
      have e0 := ih k (by omega)
      simp only [arim]
      omega

theorem dspoke_r4 (a c : ℕ → ℤ)
    (ha : ∀ k : ℕ, a (k+4) + a k = 2 * a (k+3) + 4 * a (k+2))
    (hc : ∀ k : ℕ, 3 ≤ k → c (k+4) + c k = 2 * c (k+3) + 4 * c (k+2)) :
    ∀ k : ℕ, (a (k+7) + 2 * c (k+7) + a (k+5) - a (k+6)) + (a (k+3) + 2 * c (k+3) + a (k+1) - a (k+2))
      = 2 * (a (k+6) + 2 * c (k+6) + a (k+4) - a (k+5)) + 4 * (a (k+5) + 2 * c (k+5) + a (k+3) - a (k+4)) := by
  intro k
  -- omega 原子坑规避：ha (k+j) 展开为 a (k+j+4) 等嵌套形，须复述为与目标同形的 a (k+7) 等。
  have ha1 : a (k + 5) + a (k + 1) = 2 * a (k + 4) + 4 * a (k + 3) := by
    have h := ha (k + 1)
    rw [show k + 1 + 4 = k + 5 from by omega,
        show k + 1 + 3 = k + 4 from by omega,
        show k + 1 + 2 = k + 3 from by omega] at h
    exact h
  have ha2 : a (k + 6) + a (k + 2) = 2 * a (k + 5) + 4 * a (k + 4) := by
    have h := ha (k + 2)
    rw [show k + 2 + 4 = k + 6 from by omega,
        show k + 2 + 3 = k + 5 from by omega,
        show k + 2 + 2 = k + 4 from by omega] at h
    exact h
  have ha3 : a (k + 7) + a (k + 3) = 2 * a (k + 6) + 4 * a (k + 5) := by
    have h := ha (k + 3)
    rw [show k + 3 + 4 = k + 7 from by omega,
        show k + 3 + 3 = k + 6 from by omega,
        show k + 3 + 2 = k + 5 from by omega] at h
    exact h
  have hc3 : c (k + 7) + c (k + 3) = 2 * c (k + 6) + 4 * c (k + 5) := by
    have h := hc (k + 3) (by omega)
    rw [show k + 3 + 4 = k + 7 from by omega,
        show k + 3 + 3 = k + 6 from by omega,
        show k + 3 + 2 = k + 5 from by omega] at h
    exact h
  omega

theorem drim_r4 (a c : ℕ → ℤ)
    (ha : ∀ k : ℕ, a (k+4) + a k = 2 * a (k+3) + 4 * a (k+2))
    (hc : ∀ k : ℕ, 3 ≤ k → c (k+4) + c k = 2 * c (k+3) + 4 * c (k+2)) :
    ∀ k : ℕ, (a (k+7) + c (k+7)) + (a (k+3) + c (k+3)) = 2 * (a (k+6) + c (k+6)) + 4 * (a (k+5) + c (k+5)) := by
  intro k
  have ha3 : a (k + 7) + a (k + 3) = 2 * a (k + 6) + 4 * a (k + 5) := by
    have h := ha (k + 3)
    rw [show k + 3 + 4 = k + 7 from by omega,
        show k + 3 + 3 = k + 6 from by omega,
        show k + 3 + 2 = k + 5 from by omega] at h
    exact h
  have hc3 : c (k + 7) + c (k + 3) = 2 * c (k + 6) + 4 * c (k + 5) := by
    have h := hc (k + 3) (by omega)
    rw [show k + 3 + 4 = k + 7 from by omega,
        show k + 3 + 3 = k + 6 from by omega,
        show k + 3 + 2 = k + 5 from by omega] at h
    exact h
  omega

-- ===== 主桥收口：具体序列实例化 + D_spoke/D_rim 强归纳 =====

/-- 具体序列（ℤ 化）：a(n) = 梯子匹配数，c(n) = Cst 匹配数。 -/
def az : ℕ → ℤ := fun k => ((hosoya (pathGraph k □ pathGraph 2) : ℕ) : ℤ)
def cz : ℕ → ℤ := fun k => ((matchCount (PrismR2.Cst k k) : ℕ) : ℤ)

@[simp] theorem az_val (k : ℕ) : az k = ((hosoya (pathGraph k □ pathGraph 2) : ℕ) : ℤ) := rfl
@[simp] theorem cz_val (k : ℕ) : cz k = ((matchCount (PrismR2.Cst k k) : ℕ) : ℤ) := rfl

/-- a-R4 具体实例（由 hosoya_ladder_rec4）。 -/
theorem az_r4 (k : ℕ) : az (k+4) + az k = 2 * az (k+3) + 4 * az (k+2) := by
  have h := hosoya_ladder_rec4 k
  simp only [az_val]
  push_cast
  omega

/-- c-R4 具体实例（由 c_rec4'，ambient-free）。 -/
theorem cz_r4 (k : ℕ) (hk : 3 ≤ k) : cz (k+4) + cz k = 2 * cz (k+3) + 4 * cz (k+2) := by
  have h := c_rec4' k hk
  simp only [cz_val]
  push_cast
  omega

/-- D_spoke(m) = a(m+3) + 2c(m+3) + a(m+1) − a(m+2)（d_spoke 的 ℤ 表达）。 -/
def Dspoke : ℕ → ℤ := fun m => az (m+3) + 2 * cz (m+3) + az (m+1) - az (m+2)
/-- D_rim(m) = a(m+3) + c(m+3)（d_rim 的 ℤ 表达）。 -/
def Drim : ℕ → ℤ := fun m => az (m+3) + cz (m+3)

@[simp] theorem Dspoke_val (m : ℕ) :
    Dspoke m = az (m+3) + 2 * cz (m+3) + az (m+1) - az (m+2) := rfl
@[simp] theorem Drim_val (m : ℕ) : Drim m = az (m+3) + cz (m+3) := rfl

theorem Dspoke_r4 (k : ℕ) : Dspoke (k+4) + Dspoke k = 2 * Dspoke (k+3) + 4 * Dspoke (k+2) := by
  simpa only [Dspoke_val, Nat.add_assoc, Nat.reduceAdd] using
    dspoke_r4 az cz az_r4 (fun j hj => cz_r4 j hj) k

theorem Drim_r4 (k : ℕ) : Drim (k+4) + Drim k = 2 * Drim (k+3) + 4 * Drim (k+2) := by
  simpa only [Drim_val, Nat.add_assoc, Nat.reduceAdd] using drim_r4 az cz az_r4 (fun j hj => cz_r4 j hj) k

-- 具体序列数值（az1..az6 / cz3..cz6）
theorem az1 : az 1 = 2 := by rw [az_val, hosoya_lad1]; norm_num
theorem az2 : az 2 = 7 := by rw [az_val, hosoya_lad2]; norm_num
theorem az3 : az 3 = 22 := by rw [az_val, hosoya_lad3]; norm_num
theorem az4 : az 4 = 71 := by rw [az_val, a4]; norm_num
theorem az5 : az 5 = 228 := by rw [az_val, a5]; norm_num
theorem az6 : az 6 = 733 := by rw [az_val, a6]; norm_num
theorem cz3 : cz 3 = 4 := by rw [cz_val, c3]; norm_num
theorem cz4 : cz 4 = 15 := by rw [cz_val, c4]; norm_num
theorem cz5 : cz 5 = 46 := by rw [cz_val, c5]; norm_num
theorem cz6 : cz 6 = 150 := by rw [cz_val, c6]; norm_num

/-- D_spoke 初值：25, 86, 271, 876。 -/
theorem Dspoke0 : Dspoke 0 = 25 := by
  simp only [Dspoke_val, Nat.zero_add, az3, cz3, az1, az2]
  norm_num
theorem Dspoke1 : Dspoke 1 = 86 := by
  simp only [Dspoke_val, Nat.reduceAdd, az4, cz4, az2, az3]
  norm_num
theorem Dspoke2 : Dspoke 2 = 271 := by
  simp only [Dspoke_val, Nat.reduceAdd, az5, cz5, az3, az4]
  norm_num
theorem Dspoke3 : Dspoke 3 = 876 := by
  simp only [Dspoke_val, Nat.reduceAdd, az6, cz6, az4, az5]
  norm_num

/-- D_rim 初值：26, 86, 274, 883。 -/
theorem Drim0 : Drim 0 = 26 := by
  simp only [Drim_val, Nat.zero_add, az3, cz3]
  norm_num
theorem Drim1 : Drim 1 = 86 := by
  simp only [Drim_val, Nat.reduceAdd, az4, cz4]
  norm_num
theorem Drim2 : Drim 2 = 274 := by
  simp only [Drim_val, Nat.reduceAdd, az5, cz5]
  norm_num
theorem Drim3 : Drim 3 = 883 := by
  simp only [Drim_val, Nat.reduceAdd, az6, cz6]
  norm_num

/-- **主桥（spoke 侧）**：ℤ 口径下 hosoya (prismDelSpoke (m+3)) = D_spoke(m)。
三式 omega 消元：hosoya_prism_mc + delspoke_eq + p_decomp。 -/
theorem dspoke_cast (m : ℕ) (hm : 3 ≤ m + 3) :
    ((hosoya (prismDelSpoke (m+3) hm) : ℕ) : ℤ) = Dspoke m := by
  have h0 := hosoya_prism_mc m
  have h1 := delspoke_eq m hm
  have h2 := p_decomp m
  simp only [Dspoke_val, az_val, cz_val]
  push_cast
  omega

/-- **主桥（rim 侧）**：ℤ 口径下 hosoya (prismDelRim (m+3)) = D_rim(m)。 -/
theorem drim_cast (m : ℕ) (hm : 3 ≤ m + 3) :
    ((hosoya (prismDelRim (m+3) hm) : ℕ) : ℤ) = Drim m := by
  have h := delrim_eq m hm
  simp only [Drim_val, az_val, cz_val]
  push_cast
  omega

/-- 强归纳收口（spoke 侧）：D_spoke(m) = aspoke(m)。 -/
theorem dspoke_val (m : ℕ) (hm : 3 ≤ m + 3) :
    ((hosoya (prismDelSpoke (m+3) hm) : ℕ) : ℤ) = aspoke m := by
  have h0 := dspoke_cast m hm
  have h1 := aspoke_strong Dspoke Dspoke0 Dspoke1 Dspoke2 Dspoke3 Dspoke_r4 m
  omega

/-- 强归纳收口（rim 侧）：D_rim(m) = arim(m)。 -/
theorem drim_val (m : ℕ) (hm : 3 ≤ m + 3) :
    ((hosoya (prismDelRim (m+3) hm) : ℕ) : ℤ) = arim m := by
  have h0 := drim_cast m hm
  have h1 := arim_strong Drim Drim0 Drim1 Drim2 Drim3 Drim_r4 m
  omega

end PrismR3

-- ===== 冻结定理头（statement.lean L90–98 逐字，占位符已替换为实证证明体） =====

/-- **主定理一（spoke 侧桥定理）**：对一切 `n ≥ 3`，删一条 spoke 的棱柱的匹配数等于序列
`aspoke (n - 3)`。其中 `aspoke` 为已冻结序列（初值 25, 86, 271, 876，四阶递推
`a(n+4) = 2·a(n+3) + 4·a(n+2) − a(n)`，下标 0 对应图侧 n = 3）。 -/
theorem hosoya_prism_del_spoke : ∀ (n : ℕ) (hn : 3 ≤ n),
    (hosoya (prismDelSpoke n hn) : ℤ) = aspoke (n - 3) := by
  intro n hn
  obtain ⟨m, rfl⟩ : ∃ m, n = m + 3 := ⟨n - 3, by omega⟩
  rw [show m + 3 - 3 = m from by omega]
  exact PrismR3.dspoke_val m hn

/-- **主定理二（rim 侧桥定理）**：对一切 `n ≥ 3`，删一条 rim 边的棱柱的匹配数等于序列
`arim (n - 3)`。其中 `arim` 为已冻结序列（初值 26, 86, 274, 883，同四阶递推）。 -/
theorem hosoya_prism_del_rim : ∀ (n : ℕ) (hn : 3 ≤ n),
    (hosoya (prismDelRim n hn) : ℤ) = arim (n - 3) := by
  intro n hn
  obtain ⟨m, rfl⟩ : ∃ m, n = m + 3 := ⟨n - 3, by omega⟩
  rw [show m + 3 - 3 = m from by omega]
  exact PrismR3.drim_val m hn

\end{lstlisting}

\section{Verification evidence}\label{sec:verif}

\begin{itemize}
\item Toolchain: Lean \texttt{v4.34.0} (\texttt{leanprover/lean4:v4.34.0}),
mathlib4 \texttt{v4.34.0}, revision \texttt{5ed29652}. Reproduction: with that
pin, a strict compile of \texttt{proofs/bridge.lean} exits with code $0$ and
no error; the remaining diagnostics are linter-style warnings only (unused
tactic), which the pipeline's discipline does not treat as grounds for
rejection.
\item Statement integrity: the frozen statement file (124 lines) has sha256
\begin{center}
\texttt{d186ff1651a4383e7e213b42ba4ec8c01}\allowbreak
\texttt{71d9e7004231db2764370cd967ca39e},
\end{center}
and the final proof file (1857 lines) has sha256
\begin{center}
\texttt{399ee81fe363f639811a3408397b14259}\allowbreak
\texttt{3ea1e5d195f3e916fd7b7c07af6c1}.
\end{center}
The two frozen theorem heads inside the proof file are byte-identical to
lines 90--91 and 96--97 of the frozen statement file (programmatic
comparison, recorded in the source task's final adjudication); the base
definitions at lines 28--47 of the proof file match the frozen snapshot's
lines 64--83 verbatim. The statement excerpt in Listing~\ref{lst:stmt} is an
ordered subsequence of the frozen file with the three declared exclusions and
no other deviation (verified mechanically by the package-local copy of
\texttt{inspect-tex.py} under \texttt{audit/}; record in
\texttt{audit/listing-verify.txt}).
\item Kernel check: strict compilation exits 0 with no \texttt{sorry} and no
\texttt{admit} (whole-file scan including comments: zero hits).
\texttt{\#print axioms} output, verbatim (11 declarations, independent
re-print of 2026-10-01; the narrower monospace is typesetting only):
\begin{lstlisting}[basicstyle=\ttfamily\footnotesize,frame=none]
'hosoya_prism_del_spoke' depends on axioms: [propext, Classical.choice, Quot.sound]
'hosoya_prism_del_rim' depends on axioms: [propext, Classical.choice, Quot.sound]
'PrismR3.delspoke_eq' depends on axioms: [propext, Classical.choice, Quot.sound]
'PrismR3.delrim_eq' depends on axioms: [propext, Classical.choice, Quot.sound]
'PrismR3.cell_C4' depends on axioms: [propext, Classical.choice, Quot.sound]
'PrismR3.c5' depends on axioms: [propext, Classical.choice, Quot.sound]
'PrismR3.c6' depends on axioms: [propext, Classical.choice, Quot.sound]
'PrismR3.dspoke_val' depends on axioms: [propext, Classical.choice, Quot.sound]
'PrismR3.drim_val' depends on axioms: [propext, Classical.choice, Quot.sound]
'PrismR3.aspoke_strong' depends on axioms: [propext, Quot.sound]
'PrismR3.arim_strong' depends on axioms: [propext, Quot.sound]
\end{lstlisting}
Every line is a subset of the whitelist \{propext, \texttt{Quot.sound},
\texttt{Classical.choice}\}; \texttt{sorryAx} does not occur, and no
\texttt{Lean.ofReduceBool} occurs (the proof contains no
\texttt{native\_decide}; the four \texttt{native\_decide} data anchors of the
frozen statement file are verification-only anchors outside the deposited
proof, exactly as the frozen file's own header notes).
\item Grading by the audit department: both frozen theorems are
\emph{complete proof} (final adjudication of 2026-10-01, seven checks all
passed: independent strict recompile, independent axiom audit over 11
declarations, freeze check, redeclaration/pollution check, non-degeneracy
spot checks, library byte-identity, boundary check; the adjudication text
ships as \texttt{audit/final-phase2-r3.txt} in this deposit). No weakening,
no added hypotheses, no unproven residue.
\item Audit chain (originals in this deposit under \texttt{audit/} unless
noted): the freeze-time sha256 record and the byte-wise symmetric-difference
conclusion are carried by the source task's audit trail in the working
repository; this deposit carries the final adjudication copy, the axiom
re-print, the listing byte-exactness record, and this note's claim check
\texttt{claim-check.md} --- added to \texttt{audit/} after the audit
department returns its adjudication, and a prerequisite for publication.
\end{itemize}

\section{Novelty evidence}\label{sec:novelty}

Adjudication copied from the pipeline's exclusivity review (C9 gate): verdict
\textbf{no occupying record found (reachable-channel scope)}. At the sequence
level the review covers 41 ledgered queries over six forms of the two slices
plus alias keywords (screening dossier, 58 raw responses), and the graph-side
increment adds 31 queries over six channels with five positive controls, all
passing; the strongest new sequence-side evidence is the exhaustive sweep of
the OEIS linear-recurrence index page for the signature $(2,4,0,-1)$: exactly
four sequences are registered under that signature (A102080 the mother prism
sequence, A020877 the M\"obius-ladder matchings, and two tiling/array
families), and neither slice is among them. At the graph-side proposition
level (prism with one edge deleted, endpoints kept, matching count satisfying
the fourth-order recurrence) the arXiv / Crossref / MSE / GitHub / local
mathlib channels returned no same-form statement. Value tier: \textbf{new
sequence} --- the ``new sequence'' half holds (both slices unregistered), the
``new recurrence'' half does \emph{not} hold and is not claimed: the
recurrence itself is the mother sequence A102080's registered formula (same
signature, different initial values; the review's
same-spectrum-different-initial-values wording, copied verbatim). The
graph-side semantic bridge is graded \textbf{first formalisation (Lean
three-layer zero same-form)}. The full query log, including every zero-hit
query, is in this deposit at \texttt{audit/c9-record.md} (verbatim,
byte-for-byte, three source records concatenated with their headers and
adjudication conclusions unchanged).

\begin{center}
\resizebox{\textwidth}{!}{%
\begin{tabular}{lll}
\hline
channel & queries & result \\
\hline
OEIS, numeric (screening) & six forms of the two slices (main,
drop-first variants, interleavings, subcolumns) & all zero hits; mother
A102080 confirmed in register \cite{a102080} \\
OEIS, signature index & signature $(2,4,0,-1)$ index page & 4 registered
entries, neither slice among them; strongest sequence-side closure \\
OEIS, keyword/alias & alias and keyword strings & zero this-object hits \\
arXiv & graph-side 6 + screening layer & nearest neighbours read at
abstract/full-text level (perfect-matching structural properties, PH-property,
keyword collisions): no matching count, no edge deletion \\
Crossref & graph-side 4 & one family-relevant record (Aldred--Plummer,
matching extension in prism graphs \cite{aldred2017}): structural
$n$-extendability, not a counting result; full text behind paywall, recorded
as unread \\
Mathematics Stack Exchange & graph-side 4 + control & control hit; slice
queries zero \\
GitHub Lean ecosystem & 8 & zero same-form; prism/hosoya/matching-count
intersections all empty; one prism-related Lean formalisation found and
excluded on different form (domination number, not matching count) \\
local mathlib (pinned v4.34.0) & 3 & no prism definition, no Hosoya word
family, matching predicates only (no counting API) \\
residual exposure & --- & journal full-text layer (paywalled) and general web
search not reachable; recorded as unverified, not as evidence; Aldred--Plummer
and the ladder-side Farrell 1983 record (object = ladders, not prisms) remain
unread \\
\hline
\end{tabular}%
}
\end{center}

The channels are named so a later reader can re-run them; if any channel later
returns a same-form record, the tier above weakens accordingly, and this
deposit will be amended by a later version rather than by editing this record.
A zero hit is not novelty; the strongest claim licensed here is that no record
of this object was found in the channels queried.

\section{Scope and limitations}\label{sec:limits}

Grade and boundary, stated as the audit states them. (a) The kernel layer
proves Theorems~\ref{thm:spoke} and~\ref{thm:rim} and nothing else beyond the
supporting declarations contained in the proof file. (b) The grading of both
is \emph{complete proof} ($0$ \texttt{sorry}, axiom whitelist, no scaffold
residue); this note makes no claim about recurrence minimality, closed forms,
asymptotics, or the mod-$k$ residue patterns of the sequences (the congruence
layer of the parent program is not part of this deposit). (c) The
identification of the counting function with the Hosoya index is definitional
in the formalisation (\texttt{hosoya} is defined as that edge-subset count);
no further combinatorial conjecture layer is carried. (d) The four
\texttt{native\_decide} data anchors of the frozen statement file
($n=3,4$ values $25, 86, 26, 86$) are verification-only anchors and are not
part of the deposited proof or its axiom set. (e) The novelty record's
literature layer excludes paywalled full texts; the recorded exposures
(Aldred--Plummer 2017 full text; the ladder-side Farrell 1983 record, whose
object is ladders not prisms) are recorded as unread rather than as evidence,
and no universal claim of the form ``the literature has not studied this'' is
made. (f) This note contains no literature review; it is a deposit record,
not a survey paper.

\section*{Data availability}

Lean sources, verification and novelty-search evidence are deposited on
Zenodo; the DOI is reserved at upload by the author and is recorded in
\texttt{metadata/README.md} of this package once created; see
\texttt{UPLOAD.md} in the source repository. The deposit contains the claim
note (LaTeX source and PDF), the Lean sources (\texttt{proofs/}: frozen
statement file, final proof file, \texttt{lean-toolchain}), and the audit and
novelty-search records (\texttt{audit/}). Claim note under CC BY 4.0, code
under MIT.

\section*{Use of generative AI}

\begin{itemize}
\item \emph{AI-performed stages}: candidate generation and the exclusivity
review, statement formalisation and freezing, the semantic roundtrip check,
three-phase proof search (worker agents plus main-agent scaffold assembly
under the pipeline's verified-scaffold clause, provenance recorded in the
proof file header and the task ledger), and the assembly of this note
including its verbatim listings.
\item \emph{Model, node, budget}: the proving and audit phases ran on the
session node (wire designation \texttt{glm-5.3} as recorded in the task
ledger's node-drift incident line), with no undeclared cross-node switching;
consumed: agent-run samples $12$ of the task-level permitted $20$ (Phase 0: 2,
Phase 1: 6, Phase 2: 12 -- of which the two bridge theorems are the Phase 2
r3 close-out), formal kernel compiles $17$ of the permitted $36$,
llm-calls $4$ of $8$ (plus $4$ transmission failures, ledgered separately),
and compile-debug probes $23$ (single-entry category, no cap); LaTeX $0$ on
the source task. The numbers are copied from the source task's
\texttt{budget.log} (tally line of 2026-10-01, categories as ledgered there).
This note consumed $0$ pipeline LLM calls against a cap of $4$; drafting was
mechanical assembly.
\item \emph{Final adjudication}: all statements and proofs reported here were
checked by the Lean~4 kernel: every proof compiles with no \texttt{sorry}, and
\texttt{\#print axioms} reports only \{propext, \texttt{Quot.sound},
\texttt{Classical.choice}\}. Statement text was frozen before proving and
compared byte-wise afterwards.
\item The author reviewed the material and is responsible for all content. AI
systems are not listed as authors.
\end{itemize}

\section*{Author contributions}
The author determined the research question and the scope fences, verified
the semantics of the statements, and is responsible for signing off and for
any deposit or release. The automated pipeline (LLM-driven) generated the
candidate propositions, the formalisation, the proof search and the draft
text. No external release is performed by the pipeline.

\begin{thebibliography}{9}
% Only works actually cited; each was checked to exist before entry (2026-10-01:
% lean4/mathlib via Crossref title metadata; a102080 via live oeis.org fetch;
% aldred2017 via live Crossref record). No entry is quoted from memory.
\bibitem{lean4}
L.~de Moura and S.~Ullrich.
\newblock The Lean 4 theorem prover and programming language.
\newblock \emph{Automated Deduction -- CADE-28}, LNCS, pp.~625--635, 2021.
\url{https://doi.org/10.1007/978-3-030-79876-5_37}

\bibitem{mathlib}
The mathlib Community.
\newblock The Lean mathematical library.
\newblock \emph{CPP 2020}, pp.~367--381, 2020.
\url{https://doi.org/10.1145/3372885.3373824}

\bibitem{a102080}
OEIS Foundation.
\newblock Number of matchings in the $C_n \times P_2$ ($n$-prism) graph
(sequence A102080; the mother sequence, registered under the same fourth-order
recurrence with different initial values).
\newblock \emph{The On-Line Encyclopedia of Integer Sequences}, entry
A102080.
\url{https://oeis.org/A102080}

\bibitem{aldred2017}
R.~E.~L. Aldred and M.~D. Plummer.
\newblock Matching extension in prism graphs.
\newblock \emph{Discrete Applied Mathematics} 221 (2017), pp.~25--32.
\url{https://doi.org/10.1016/j.dam.2016.12.017}
\end{thebibliography}

\end{document}
